
\documentclass[11pt]{article}

\usepackage[a4paper,margin=1in]{geometry}
\usepackage{fontspec}
\usepackage{microtype}
\DisableLigatures{encoding = *, family = *}
\usepackage{amsmath,amssymb,amsthm,mathtools}
\usepackage{booktabs,longtable,array}
\usepackage{calc}
\usepackage{enumitem}
\usepackage{xurl}
\usepackage[hidelinks]{hyperref}
\hypersetup{
  pdftitle={Finite Sigma-Equivariant Row-Graded Magmas on AG(2,3): Landscape, Intrinsic Selection, and Linearization},
  pdfauthor={Burundai Taryi},
  pdfsubject={Finite sigma-equivariant row-graded magmas on AG(2,3)},
  pdfkeywords={finite magma, row-graded magma, affine plane, AG(2,3), sigma-equivariance, finite field, associativity atlas, intrinsic selection, linearization}
}
\usepackage{bookmark}

\setlength{\parskip}{0.45em}
\setlength{\parindent}{0pt}
\setcounter{secnumdepth}{3}
\setcounter{tocdepth}{2}
\setlength{\LTleft}{0pt}
\setlength{\LTright}{0pt}
\sloppy
\emergencystretch=3em
\providecommand{\tightlist}{%
  \setlength{\itemsep}{0pt}\setlength{\parskip}{0pt}}

\theoremstyle{definition}
\newtheorem{theorem}{Theorem}[section]
\newtheorem{lemma}[theorem]{Lemma}
\newtheorem{proposition}[theorem]{Proposition}
\newtheorem{corollary}[theorem]{Corollary}
\newtheorem*{mainthm}{Main theorem package}
\renewcommand{\qedsymbol}{$\square$}
\newcommand{\RP}{\mathsf{RP}}
\newcommand{\RC}{\mathsf{RC}}

\makeatletter
\renewcommand{\verbatim@font}{\normalfont\ttfamily\small}
\makeatother




\title{Finite \(\sigma\)-Equivariant Row-Graded Magmas on \(\operatorname{AG}(2,3)\): Landscape, Intrinsic Selection, and Linearization}
\author{Burundai Taryi\thanks{Independent researcher, \href{mailto:burundai.taryi@gmail.com}{burundai.taryi@gmail.com}.}}
\date{}


\begin{document}
\maketitle

\begin{abstract}
We study the finite landscape of \(\sigma\)-equivariant row-graded magmas on
\(M=S\times S\), \(S=\mathbb F_3\), equivalently the point set of
\(\operatorname{AG}(2,3)\), with left-row output and a fixed same-row
off-diagonal Steiner rule.  The landscape has normal form
\[
\Omega'\cong S^{18}\times S^3\cong S^{21},
\qquad |\Omega'|=3^{21}=10\,460\,353\,203.
\]

Two finite atlas theorems first map the ambient space.  The associativity atlas
satisfies
\[
63\le \operatorname{Assoc}\le 597,
\]
with exact endpoint loci and partition function.  The hidden-continuation
frontier satisfies
\[
\min H_{tot}=-2268,
\qquad
\max H_{tot}=7302,
\qquad
\max_{N_-=0}H_{tot}=7020.
\]
Inside this landscape, four internal finite predicates
\[
\mathcal H_{acc}=0,
\qquad
H_{diag}=0,
\qquad
(r,r)^2=(r,r),
\qquad
\operatorname{AbsTrans}=\{C,J\}
\]
select a unique point, the right-row projection magma \(\RP\).  After
selection, \(\RP\) has a finite linearization passport: rank-three left
translations, a \(51\)-element positive composition envelope, an endomorphism
and intertwiner classification, a binary term algebra of size \(630\), a
row-contraction matrix bridge, a hidden-continuation operator identity,
carrier-level compact companions, and a signed symplectic lift of the pure
\(C/J\) geometry.

The theorem package is organized as
\[
\text{finite landscape}
\Longrightarrow
\text{intrinsic finite selection}
\Longrightarrow
\text{linearization of }\RP.
\]
\end{abstract}

\tableofcontents
\newpage

\phantomsection
\section*{Introduction and main theorem}
\addcontentsline{toc}{section}{Introduction and main theorem}
\label{introduction-main-theorem}


\subsection*{Object}\label{intro-object}

Let
\[
S=\mathbb F_3,
\qquad
M=S\times S.
\]
The first coordinate is the row and the second coordinate is the column.  The
paper studies binary operations on \(M\) satisfying left-row output, the fixed
same-row off-diagonal Steiner rule, and equivariance under the simultaneous
shift \(\sigma(r,c)=(r+1,c+1)\).  The normal coordinates are
\[
(A,B,d)\in S^9\times S^9\times S^3,
\]
where \(A\) and \(B\) encode the two ordered cross-row representatives and
\(d\) encodes the same-point diagonal.  The resulting finite landscape is
\[
\Omega'\cong S^{18}\times S^3\cong S^{21}.
\]

\subsection*{Main theorem package}\label{intro-main-theorem-package}

\begin{mainthm}
Let \(\Omega'\) be the row-graded landscape on \(M=S\times S\) satisfying
left-row output, the fixed same-row off-diagonal Steiner rule, and
\(\sigma\)-equivariance.  Then:
\begin{enumerate}[leftmargin=2em]
\item The landscape has normal form
\[
\Omega'\cong S^{18}\times S^3\cong S^{21},
\qquad
|\Omega'|=3^{21}=10\,460\,353\,203.
\]
\item The associativity atlas satisfies
\[
63\le\operatorname{Assoc}\le597,
\]
with exact endpoint loci and partition function.  The hidden-continuation
frontier satisfies
\[
\min H_{tot}=-2268,
\qquad
\max H_{tot}=7302,
\qquad
\max_{N_-=0}H_{tot}=7020.
\]
\item The four internal selector predicates
\[
\mathcal H_{acc}=0,
\qquad
H_{diag}=0,
\qquad
(r,r)^2=(r,r),
\qquad
\operatorname{AbsTrans}=\{C,J\}
\]
select the right-row projection magma \(\RP\) uniquely, with multiplication
\[
(r_1,c_1)\cdot(r_2,c_2)=
\begin{cases}
(r_1,r_2), & r_1\ne r_2,\\[2mm]
(r_1,\overline{c_1c_2}), & r_1=r_2,\ c_1\ne c_2,\\[2mm]
(r_1,r_1), & r_1=r_2,\ c_1=c_2.
\end{cases}
\]
\item The selected magma has the post-selection passport proved in Part III:
left translations and positive composition envelope, endomorphisms and
intertwiners, binary term functions, matrix row-contraction, a
hidden-continuation operator bridge, carrier-level compact companions, and a
signed symplectic lift.
\item The finite atlas and passport claims are supported by the certificate
protocol recorded in the companion material.
\end{enumerate}
\end{mainthm}

\begin{proof}
The normal form is Theorem~\ref{thm:1.1}.  The associativity atlas is
Theorem~\ref{thm:5.2}, with endpoint and partition data in
Theorems~\ref{thm:6.1} and~\ref{thm:6.2}.  The hidden-continuation frontier is
Theorem~\ref{thm:7.2}.  The selection statement is
Theorem~\ref{thm:12.3}.  The post-selection passport is given by
Lemma~\ref{lem:13.1}, Theorems~\ref{thm:14.2}, \ref{thm:14.3},
\ref{thm:15.2}, \ref{thm:15.4}, \ref{thm:16.2}, and~\ref{thm:18.3},
together with Propositions~\ref{prop:15.3},
\ref{prop:constant-intertwiners}, \ref{prop:intertwiner-classification},
and~\ref{prop:16.3}, and Theorems~\ref{thm:17.1}, \ref{thm:19.1},
and~\ref{thm:20.1}.  The certificate interface is recorded in
Section~\ref{finite-certificate-protocol}.
\end{proof}

\subsection*{Proof architecture}\label{intro-proof-architecture}

The proof separates the roles of the invariants and predicates.
\begin{longtable}{@{}p{0.27\linewidth}p{0.65\linewidth}@{}}
\toprule
component & role \\
\midrule
\endhead
\bottomrule
\endlastfoot
Part I: finite landscape & normal form, effective symmetry, associativity atlas, and hidden-continuation frontier \\
Part II: intrinsic selection & four selector predicates and the uniqueness proof for \(\RP\) \\
Part III: linearization & algebraic and carrier-level structures exposed after \(\RP\) is fixed \\
Companion material & finite certificate protocol, \(\RP\) one-page passport, structural reading, open bridges, and assistance note \\
\end{longtable}
Associativity and the \(\mathcal H\)-frontier are atlas invariants.  The four
conditions in the main theorem package are selector predicates.  The compact Lie
groups appearing in Part III are carrier-level compact companions attached to
specified linear structures on the shared carrier \(k[M]\cong\operatorname{Mat}_3(k)\).

\subsection*{Guide to the paper}\label{intro-guide}

Part I constructs the finite landscape and records its two global atlases.  Part
II applies the internal selector and proves uniqueness of \(\RP\).  Part III
fixes the selected multiplication and develops its operator, term-function,
endomorphism, matrix, compact-companion, and symplectic structures.  The
companion material records the certificate path, a compact passport for the
selected point \(\RP\), the row-complement comparison, the Landauer/Shannon
structural reading, and open bridge problems.

\part{The finite landscape}

Part I constructs the ambient finite space and records the two global atlases used later: associativity and hidden continuation.  The statements here describe the landscape before any selector is applied.

\section{Carrier, axioms, and normal form}\label{carrier-axioms-normal-form}\label{carrier-axioms-and-normal-form}

\subsection{Carrier and shift}\label{carrier-and-shift}

Let
\[
S=\mathbb F_3,
\qquad
M=S\times S.
\]
All arithmetic in row and column coordinates is in \(S\).  Equivalently,
\(M=\mathbb F_3^2\) is the point set of the affine plane
\(\operatorname{AG}(2,3)\), with the chosen row--column coordinates fixed
throughout.  For \(x=(r,c)\in M\), write
\[
\operatorname{row}(x)=r,
\qquad
\operatorname{col}(x)=c,
\]
and let
\[
F_r=\{r\}\times S
\]
be the row fiber.  We use the cyclic shift
\[
\sigma(r,c)=(r+1,c+1).
\]
Thus \(\sigma\) translates rows and columns simultaneously.

The row-relative offset \(c-r\) is invariant under \(\sigma\).  Hence \(M\)
decomposes into the three \(\sigma\)-orbits
\[
D_i=\{(r,r+i):r\in S\},
\qquad i\in S.
\]
The orbit \(D_0\) is the diagonal orbit, while \(D_1\) and \(D_2\) are the two
non-diagonal orbits.

For distinct \(i,j\in S\), write \(\overline{ij}\) for the third element
of \(S\setminus\{i,j\}\).  Equivalently, in \(\mathbb F_3\),
\[
\overline{ij}=-i-j
\qquad (i\ne j).
\]

\subsection{Axioms}\label{axioms}

The landscape \(\Omega'\) consists of binary operations \(\cdot\) on \(M\)
satisfying the following three conditions.

\begin{description}[leftmargin=3.2em,style=nextline]
\item[AX0, left-row output.]
For all \(x,y\in M\),
\[
\operatorname{row}(x\cdot y)=\operatorname{row}(x).
\]
Equivalently, \(x\cdot y\in F_{\operatorname{row}(x)}\).

\item[\(\sigma\)-equivariance.]
For all \(x,y\in M\),
\[
\sigma(x\cdot y)=\sigma(x)\cdot \sigma(y).
\]

\item[AX2, same-row off-diagonal Steiner rule.]
For all \(r,i,j\in S\) with \(i\ne j\),
\[
(r,r+i)\cdot(r,r+j)=(r,r+\overline{ij}).
\]
Equivalently, in absolute column notation,
\[
(r,c_1)\cdot(r,c_2)=(r,\overline{c_1c_2})
\qquad(c_1\ne c_2).
\]
This is the Steiner multiplication on a three-point row fiber
\cite{Pflugfelder1990,ColbournDinitz2007}.
\end{description}

AX0 fixes only the output row.  The output column remains free in the
cross-row branch and on the same-point diagonal branch.

\subsection{Normal coordinates}\label{normal-coordinates}

It is convenient to use row-relative coordinates adapted to this orbit
decomposition.  A point of row \(r\) in \(D_i\) is written as \((r,r+i)\),
where \(i\in S\) is the column offset from the row.  Thus \(i\) is the
\(\sigma\)-orbit label.

A normal coordinate is a triple
\[
(A,B,d)\in S^9\times S^9\times S^3,
\]
where
\[
A=(A_{ij})_{i,j\in S},
\qquad
B=(B_{ij})_{i,j\in S},
\qquad
 d=(d_0,d_1,d_2).
\]
The associated multiplication is defined by the four formulas
\[
(r,r+i)\cdot(r+1,r+1+j)=(r,r+A_{ij}),
\]
\[
(r,r+i)\cdot(r+2,r+2+j)=(r,r+B_{ij}),
\]
\[
(r,r+i)\cdot(r,r+i)=(r,r+d_i),
\]
and, for \(i\ne j\),
\[
(r,r+i)\cdot(r,r+j)=(r,r+\overline{ij}).
\]
The first two arrays record the two ordered cross-row representatives
\((0,1)\) and \((0,2)\).  The vector \(d\) records the same-point branch by
\(\sigma\)-orbit label: \(d_i\) is the output offset of \(x^2\) for
\(x\in D_i\).  The fourth formula is AX2.

Conversely, any operation satisfying AX0, \(\sigma\)-equivariance, and AX2
has a unique coordinate triple \((A,B,d)\): after applying a power of
\(\sigma\), every cross-row product reduces to one of the two row-pair
representatives above, and every same-point product reduces to one of the
three diagonal coordinates.

\subsection{Normal form theorem}\label{normal-form-theorem}

\begin{theorem}
\label{thm:1.1}
The AX0+AX2, \(\sigma\)-equivariant landscape has normal form
\[
\Omega'\cong S^{18}\times S^3\cong S^{21}.
\]
Consequently
\[
|\Omega'|=3^{21}=10\,460\,353\,203.
\]
\end{theorem}

\begin{proof}
AX0 fixes the output row of every product.  AX2 fixes all same-row
off-diagonal output columns.  The same-point diagonal branch has three
row-relative coordinates \(d_0,d_1,d_2\).  For cross-row products,
\(\sigma\)-equivariance reduces the ordered row-pair data to the two
representatives \((0,1)\) and \((0,2)\).  Each representative has
\(3\times3\) column-offset inputs, hence contributes nine independent
coordinates.  Thus the cross-row branch contributes \(18\) coordinates and
the diagonal branch contributes \(3\) coordinates.

The coordinate construction above gives a product satisfying AX0,
\(\sigma\)-equivariance, and AX2 for every \((A,B,d)\in S^9\times S^9
\times S^3\), and the preceding paragraph recovers these coordinates
uniquely from any such product.  Hence \(\Omega'\cong S^{21}\), and its
cardinality is \(3^{21}\).
\end{proof}


\section{Effective symmetry}\label{effective-symmetry}\label{effective-symmetry-and-orbit-count}

The oriented row-graded setup fixes the generator
\(s\mapsto s+1\) through
\(\sigma(r,c)=(r+1,c+1)\).  For landscape compression we use the normalizer of this cyclic shift in
\(\operatorname{Sym}(S)\).  Translations
\(s\mapsto s+\beta\) commute with the shift and act trivially on the
row-relative normal coordinates.  The remaining normalizer class is
represented by the reflection \(\nu(s)=-s\), which conjugates
\(\sigma\) to \(\sigma^{-1}\).  Since \(\sigma\)-equivariance is equivalent
to \(\sigma^{-1}\)-equivariance, this reflection preserves \(\Omega'\) as a
landscape-compression symmetry.

On normal coordinates it induces the involution
\[
\iota(A,B,d)=(A',B',d')
\]
with indices in \(S\) and
\[
A'_{ij}=-B_{-i,-j},
\qquad
B'_{ij}=-A_{-i,-j},
\qquad
d'_i=-d_{-i}.
\]
We denote the resulting residual effective action by
\[
G_{\mathrm{eff}}=\langle\iota\rangle\cong C_2.
\]
On the diagonal coordinate sector \(\Delta=S^3\), the same residual
reflection gives the compression
\[
|\Delta/G_{\mathrm{eff}}|=\frac{3^3+3}{2}=15.
\]

\begin{proposition}
\label{prop:2.1}
The effective orbit count of the full
normal-coordinate landscape is
\[
|\Omega'/G_{\mathrm{eff}}|=\frac{3^{21}+3^{10}}2.
\]
\end{proposition}

\begin{proof}
Burnside's lemma applies to the displayed involution.  The identity
contributes \(3^{21}\).  For a fixed point of \(\iota\), the nine entries
of \(A\) are free and then determine \(B\) by
\(B_{ij}=-A_{-i,-j}\).  The diagonal equations are
\[
d_0=-d_0,
\qquad
d_1=-d_2,
\qquad
d_2=-d_1,
\]
so \(d_0=0\) and \(d_1\) is free.  Hence the fixed locus has
\(9+1=10\) ternary degrees of freedom and size \(3^{10}\).  Therefore
\[
|\Omega'/G_{\mathrm{eff}}|=
\frac{3^{21}+3^{10}}2.
\]
\end{proof}

The orbit count is a landscape compression.  The residual reflection records
the finite normal-coordinate symmetry available at the atlas level.  The
directed-edge geometry used later is formulated on \(M^\times\) itself.


\section{Associativity atlas}\label{associativity-atlas}\label{associativity-master-formula}

\subsection{Block decomposition}\label{assoc-block-decomposition}\label{the-associativity-statistic}

The associativity statistic is a global atlas invariant on \(\Omega'\).  For \(x\in\Omega'\), define \[
\operatorname{Assoc}(x)=
\#\{(u,v,w)\in M^3:(u\cdot v)\cdot w=u\cdot(v\cdot w)\}.
\] Since \(|M|=9\), this statistic lies between \(0\) and \(729\).

\begin{theorem}
\label{thm:4.1}
In normal coordinates, \[
\operatorname{Assoc}=3\,rawAssoc,
\] where \(rawAssoc\) is a normalized sum over \(3^5=243\) triples and
decomposes as \[
rawAssoc=rawRRR+rawRRS+rawRSR+rawSRR+rawDIST.
\] The five blocks correspond to the row patterns \[
RRR,
\qquad RRS,
\qquad RSR,
\qquad SRR,
\qquad DIST.
\]
\end{theorem}

\begin{proof}
By \(\sigma\)-equivariance the first row coordinate of a
triple may be normalized to \(0\). The two remaining row variables and
the three column variables give \(3^5\) normalized cases; the factor
\(3\) restores the three translates. The displayed blocks are the five
equality patterns of the three rows.
\end{proof}

\subsection{Same-row block}\label{assoc-same-row-block}\label{the-same-row-block}

The all-same-row block depends only on the diagonal vector \(d\). Let \[
E(d)=\#\{(i,j):i\ne j,
\ d_i=d_j\},
\qquad
N(d)=\#\{(i,j):i\ne j,
\ d_i=j,
\ d_j=j\}.
\]

\begin{lemma}
\label{lem:4.2}
The same-row contribution is \[
rawRRR(d)=9+E(d)+2N(d).
\] The four diagonal classes give raw same-row values \[
9,\qquad 11,\qquad 13,\qquad 19.
\] In particular, \[
rawRRR(000)=19.
\]
\end{lemma}

\begin{proof}
In one row, unequal column pairs use the Steiner rule
and equal pairs use \(d\). The \(27\) triples split into a baseline
family plus the two diagonal-sensitive families counted by \(E(d)\) and
\(N(d)\). This gives the formula and the four class values.
\end{proof}

\subsection[Column-blind profile at d=000]{Column-blind profile at \(d=000\)}\label{assoc-column-blind-profile}\label{column-blind-profile-at-d000}

A cross-row rule is column-blind when its output column depends only on the two row coordinates. By \(\sigma\)-equivariance such a rule is determined by two anchors \(a=h(0,1)\) and \(b=h(0,2)\). Let \(h_{a,b}\) be the column-blind rule \[
h(r,r+1)=r+a,
\qquad
h(r,r+2)=r+b.
\] At \(d=000\), same-point products satisfy \((r,c)^2=(r,r)\).

\begin{proposition}
\label{prop:4.3}
On the column-blind, \(d=000\) stratum, \[
\begin{aligned}
&(rawRRR,rawRRS,rawRSR,rawSRR,rawDIST)(a,b,000)\\
&\quad =
\bigl(19,\, 9[a=0]+9[b=0],\, 9[a=0]+9[b=0],\\
&\hspace{7em}54,\,54[a=b]\bigr).
\end{aligned}
\] Consequently \[
\boxed{
\operatorname{Assoc}_{000}(h_{a,b})
=
3\left(73+18[a=0]+18[b=0]+54[a=b]\right).
}
\]
\end{proposition}

\begin{proof}
The \(RRR\) value is Lemma~\ref{lem:4.2}. In the \(RRS\) and
\(RSR\) patterns, the same-row step must return the cross-row anchor; at
\(d=000\) this contributes \(9\) exactly when the relevant anchor is
\(0\). In \(SRR\), the final cross-row output ignores same-row column
processing, giving all \(54\) normalized triples. In \(DIST\), the two
association orders compare the two cross-row anchors, so all \(54\)
triples survive exactly when \(a=b\).
\end{proof}

The nine column-blind rules split as follows:

\begin{longtable}[]{@{}
  >{\raggedright\arraybackslash}p{(\columnwidth - 6\tabcolsep) * \real{0.2400}}
  >{\raggedright\arraybackslash}p{(\columnwidth - 6\tabcolsep) * \real{0.3000}}
  >{\raggedleft\arraybackslash}p{(\columnwidth - 6\tabcolsep) * \real{0.2600}}
  >{\raggedleft\arraybackslash}p{(\columnwidth - 6\tabcolsep) * \real{0.2000}}@{}}
\toprule\noalign{}
\begin{minipage}[b]{\linewidth}\raggedright
class
\end{minipage} & \begin{minipage}[b]{\linewidth}\raggedright
representatives
\end{minipage} & \begin{minipage}[b]{\linewidth}\raggedleft
profile
\end{minipage} & \begin{minipage}[b]{\linewidth}\raggedleft
\(\operatorname{Assoc}_{000}\)
\end{minipage} \\
\midrule\noalign{}
\endhead
\bottomrule\noalign{}
\endlastfoot
anchored trivial & \((0,0)\) & \((19,18,18,54,54)\) & \(489\) \\
off-anchor trivial & \((1,1),(2,2)\) & \((19,0,0,54,54)\) & \(381\) \\
partly anchored nondegenerate & \((0,1),(0,2),(1,0),(2,0)\) &
\((19,9,9,54,0)\) & \(273\) \\
fully off-anchor nondegenerate & \((1,2),(2,1)\) & \((19,0,0,54,0)\) &
\(219\) \\
\end{longtable}

The fourth class consists of the selected point \(\RP=(a,b,d)=(1,2,000)\)
and the row-complement comparison point \(\RC=(a,b,d)=(2,1,000)\).

\subsection{Global range}\label{assoc-global-range}\label{global-associativity-range}\label{exact-range}

\begin{lemma}[Associativity atlas certificate protocol]
\label{lem:5.1}
The
global associativity atlas is used through five finite obligations.

\begin{enumerate}
\def\labelenumi{\arabic{enumi}.}
\tightlist
\item
  \textbf{A1, block reduction.} The associativity statistic is reduced
  by \(\sigma\)-normalization to the five row-pattern blocks \[
  rawAssoc=rawRRR+rawRRS+rawRSR+rawSRR+rawDIST
  \] of Theorem~\ref{thm:4.1}.
\item
  \textbf{A2, diagonal decomposition.} The same-row block is evaluated
  from the diagonal vector by Lemma~\ref{lem:4.2}, giving the four possible values
  \[
  9,\quad 11,\quad 13,\quad 19.
  \]
\item
  \textbf{A3, fixed-diagonal cross-range certificates.} For each
  diagonal class, the certificate verifies the exact range of \[
  rawCross=rawRRS+rawRSR+rawSRR+rawDIST.
  \]
\item
  \textbf{A4, endpoint witnesses and endpoint counts.} The endpoint
  certificates verify the normal-coordinate witnesses, endpoint counts,
  and orbit representatives for \(\operatorname{Assoc}=63\) and
  \(\operatorname{Assoc}=597\).
\item
  \textbf{A5, distribution consistency.} The partition-function
  certificate gives the full coefficient table for \[
  Z(q)=\sum_{x\in\Omega'}q^{\operatorname{Assoc}(x)}
  \] and checks total mass, range, moments, and mode.
\end{enumerate}

Any certificate satisfying A1--A5 proves the global associativity range,
endpoint loci, and partition-function claims in this atlas section.
\hfill$\square$
\end{lemma}

\begin{theorem}[Global associativity atlas]
\label{thm:5.2}
On \(\Omega'\), \[
21\le rawAssoc\le199.
\] Equivalently, \[
63\le\operatorname{Assoc}\le597.
\]

The finite certificate is organized by diagonal class:

\begin{longtable}[]{@{}
  >{\raggedright\arraybackslash}p{(\columnwidth - 6\tabcolsep) * \real{0.2000}}
  >{\raggedleft\arraybackslash}p{(\columnwidth - 6\tabcolsep) * \real{0.2667}}
  >{\raggedleft\arraybackslash}p{(\columnwidth - 6\tabcolsep) * \real{0.2667}}
  >{\raggedleft\arraybackslash}p{(\columnwidth - 6\tabcolsep) * \real{0.2667}}@{}}
\toprule\noalign{}
\begin{minipage}[b]{\linewidth}\raggedright
diagonal class
\end{minipage} & \begin{minipage}[b]{\linewidth}\raggedleft
rawRRR
\end{minipage} & \begin{minipage}[b]{\linewidth}\raggedleft
certified rawCross range
\end{minipage} & \begin{minipage}[b]{\linewidth}\raggedleft
rawAssoc consequence
\end{minipage} \\
\midrule\noalign{}
\endhead
\bottomrule\noalign{}
\endlastfoot
permutation & 9 & 12..180 & 21..189 \\
two-valued, no fixed point & 11 & 10..162 & 21..173 \\
two-valued, with fixed point & 13 & 12..174 & 25..187 \\
constant & 19 & 12..180 & 31..199 \\
\end{longtable}

Here \[
rawCross=rawRRS+rawRSR+rawSRR+rawDIST.
\] The endpoint values are therefore \(rawAssoc=21\) and
\(rawAssoc=199\), or \(\operatorname{Assoc}=63\) and
\(\operatorname{Assoc}=597\).
\end{theorem}

\begin{proof}
The hand-readable reduction is Theorem~\ref{thm:4.1} plus Lemma~\ref{lem:4.2}. After that reduction, the global claim is exactly the A3
cross-range obligation of Lemma~\ref{lem:5.1}. Adding the certified cross-row
ranges to the four possible same-row values gives the displayed global
range.

The endpoint witnesses, endpoint counts, orbit representatives, and
partition-function consistency are the A4--A5 obligations of Lemma~\ref{lem:5.1}.
They are recorded in the \texttt{L1/tables/} atlas files listed in
\texttt{SUPPORT\_INDEX.txt} and checked by
\texttt{L1/scripts/verify\_layer1\_v3\_final.py}.
\end{proof}

\paragraph{Atlas landmarks.}\label{global-and-local-associativity-roles}
The selected point has \[
\operatorname{Assoc}(\RP)=219,
\]
while the landscape minimum is \(63\).  On the column-blind \(d=000\)
profile, the value \(\operatorname{Assoc}_{000}=219\) is attained by the
two fully off-anchor nondegenerate rules \(\RP\) and \(\RC\).

\subsection{Endpoint loci and partition function}\label{assoc-endpoints-partition}\label{endpoint-loci-and-partition-function}

\subsubsection{Endpoint loci}\label{endpoint-loci}

\begin{theorem}
\label{thm:6.1}
The associativity endpoint loci have sizes \[
|\mathcal M_{63}|=24,
\qquad
|\mathcal M_{63}/\operatorname{Sym}(S)|=12,
\] and \[
|\mathcal M_{597}|=6,
\qquad
|\mathcal M_{597}/\operatorname{Sym}(S)|=3.
\] The minimum endpoint block profiles are \[
(11,2,2,2,4)
\quad\text{and}\quad
(9,2,2,4,4),
\] while the maximum endpoint block profile is \[
(19,36,36,54,54).
\]
\end{theorem}

\begin{proof}
This is the A4 endpoint part of Lemma~\ref{lem:5.1}. The endpoint
certificates record the exact solution counts by diagonal sector, orbit
representatives, and explicit witnesses in normal coordinates. The
minimum occurs in the permutation and two-valued/no-fixed diagonal
classes. The maximum occurs in the constant diagonal class. The block
profiles are obtained by evaluating the five normalized associativity
blocks on the endpoint witnesses.
\end{proof}

\subsubsection{Exact associativity partition function}\label{exact-associativity-partition-function}

\begin{theorem}
\label{thm:6.2}
The exact polynomial \[
Z(q)=\sum_{x\in\Omega'}q^{\operatorname{Assoc}(x)}
\] has \(167\) nonzero coefficients and total mass \[
Z(1)=3^{21}.
\] Its first two moments are \[
\mathbb E[\operatorname{Assoc}]=245,
\qquad
\operatorname{Var}(\operatorname{Assoc})=\frac{27820}{27}.
\] The mode occurs at \[
\operatorname{Assoc}=237
\] with coefficient \[
426,317,976.
\]
\end{theorem}

\begin{proof}
This is the A5 partition-function part of Lemma~\ref{lem:5.1}. The
coefficient certificate in \texttt{L1/tables/} is the exact finite
distribution of \(\operatorname{Assoc}\) on \(\Omega'\).  The total mass,
range, mean, variance, and mode are computed directly from that table and
checked by \texttt{L1/scripts/verify\_layer1\_v3\_final.py}.
\end{proof}


\section{Hidden-continuation frontier}\label{hidden-continuation-frontier}\label{hidden-continuation}

\subsection{Local distance reduction}\label{h-local-distance-reduction}\label{local-distance-reduction}

The hidden-continuation statistic is the second global atlas invariant.  It is a finite path-information defect.
The continuation side keeps the two-step route data, while the
left-translation side compares only the collapsed endpoint operators.
Thus \(h_q/2\) is the signed Hamming contrast between a path-resolved
comparison and an endpoint-resolved comparison. Negative local values
are precisely the local cancellation failures counted by \(N_-\). The
normalization \(H_{tot}=6rawH\) is the \(\sigma\)-orbit normalization
used by the \(\mathcal H\)-frontier certificate.

Its local index is
\[
q=(b,t,a,e,f)\in S^5.
\]
For a completed normal-form point \(x\in\Omega'\), let \(M_s\) denote
the local table read: \(M_1\) and \(M_2\) are the two cross-row tables,
and \(M_0\) is the same-row diagonal/complement rule. All row and column
arithmetic below is in \(S=\mathbb F_3\). Define
\[
xy=M_b(a,e),
\qquad
yz=M_{t-b}(e-b,f-b),
\qquad
z_R=b+yz,
\]
and
\[
ep_L=M_t(xy,f),
\qquad
ep_R=M_b(a,z_R).
\]

For \(c\in S\), define the collapsed left-translation signature
\[
L_c(i,j)=M_i(c,j)
\qquad ((i,j)\in S^2).
\]
For \(s,\alpha,z\in S\), define the stepwise continuation signature
\[
C_{s,\alpha,z}(i,j)
=
M_s\!\left(
\alpha,\,
s+M_{i-s}(z-s,j-s)
\right)
\qquad ((i,j)\in S^2).
\]
The associated Hamming distances are
\[
D_L(c,c')
=
\#\{(i,j)\in S^2:L_c(i,j)\ne L_{c'}(i,j)\},
\]
and
\[
D_C(s,\alpha,z;s',\alpha',z')
=
\#\{(i,j)\in S^2:
C_{s,\alpha,z}(i,j)\ne C_{s',\alpha',z'}(i,j)\}.
\]

\paragraph{Direct local definition.}
For each of the nine test pairs \((i,j)\in S^2\), set
\[
\delta^C_{ij}(q;x)
=
\mathbf 1\!\left\{
C_{t,xy,f}(i,j)\ne C_{b,a,z_R}(i,j)
\right\},
\]
and
\[
\delta^L_{ij}(q;x)
=
\mathbf 1\!\left\{
L_{ep_L}(i,j)\ne L_{ep_R}(i,j)
\right\}.
\]
The nine direct local summands are
\[
\Delta_{ij}(q;x)=2\bigl(\delta^C_{ij}(q;x)-\delta^L_{ij}(q;x)\bigr),
\]
displayed as
\[
\begin{array}{c|ccc}
 & j=0 & j=1 & j=2\\
\hline
i=0 & \Delta_{00} & \Delta_{01} & \Delta_{02}\\
i=1 & \Delta_{10} & \Delta_{11} & \Delta_{12}\\
i=2 & \Delta_{20} & \Delta_{21} & \Delta_{22}
\end{array}
\]
and the direct hidden-continuation contribution is
\[
h_q(x)
=
\sum_{(i,j)\in S^2}\Delta_{ij}(q;x).
\]

\begin{theorem}[Local distance reduction]
\label{thm:7.1}
For every \(x\in\Omega'\) and \(q\in S^5\),
\[
\boxed{
h_q(x)=
2\left(
D_C(t,xy,f;b,a,z_R)-D_L(ep_L,ep_R)
\right).
}
\]
Consequently the local direct definition and the Hamming-distance formula
define the same global observables \(rawH\) and \(H_{tot}\).
\end{theorem}

\begin{proof}
By definition, \(D_C(t,xy,f;b,a,z_R)\) is the sum of the nine indicators
\(\delta^C_{ij}(q;x)\), and \(D_L(ep_L,ep_R)\) is the sum of the nine
indicators \(\delta^L_{ij}(q;x)\). Subtracting these two Hamming
distances and multiplying by \(2\) gives exactly the displayed sum of
the nine direct local summands \(\Delta_{ij}\). The symbolic reduction is recorded by the \(\mathcal H\)-frontier
certificate handles \texttt{HF-RED} and \texttt{HF-CERT-IF}.
\end{proof}

Then
\[
rawH(x)=\sum_{q\in S^5}\frac{h_q(x)}2,
\qquad
H_{tot}(x)=6\,rawH(x)=3\sum_{q\in S^5}h_q(x),
\]
and
\[
\boxed{
N_-(x)=3\,\#\{q\in S^5:h_q(x)<0\}.
}
\]
Thus purity is the finite system
\[
N_-(x)=0
\quad\Longleftrightarrow\quad
h_q(x)\ge0\quad(q\in S^5).
\]
For the signed-layer diagnostics, set
\[
H_{pos}(x)=3\sum_{q\in S^5}\max\{h_q(x),0\},
\qquad
H_{\mathrm{neg\_abs}}(x)=3\sum_{q\in S^5}\max\{-h_q(x),0\}.
\]
Then
\[
H_{tot}(x)=H_{pos}(x)-H_{\mathrm{neg\_abs}}(x).
\]
Also write
\[
h_{\mathrm{loc,min}}(x)=\min_{q\in S^5}h_q(x),
\qquad
h_{\mathrm{loc,max}}(x)=\max_{q\in S^5}h_q(x).
\]

\subsection{Global frontier theorem}\label{h-global-frontier-theorem}\label{global-frontier-theorem}

The global \(\mathcal H\)-frontier theorem is supported by four finite
certificate obligations.

\begin{longtable}[]{@{}
  >{\raggedright\arraybackslash}p{(\columnwidth - 4\tabcolsep) * \real{0.3333}}
  >{\raggedright\arraybackslash}p{(\columnwidth - 4\tabcolsep) * \real{0.3333}}
  >{\raggedright\arraybackslash}p{(\columnwidth - 4\tabcolsep) * \real{0.3333}}@{}}
\toprule\noalign{}
\begin{minipage}[b]{\linewidth}\raggedright
handle
\end{minipage} & \begin{minipage}[b]{\linewidth}\raggedright
obligation
\end{minipage} & \begin{minipage}[b]{\linewidth}\raggedright
result
\end{minipage} \\
\midrule\noalign{}
\endhead
\bottomrule\noalign{}
\endlastfoot
H1 & unrestricted full-landscape range & \(rawH\in[-378,1217]\),
endpoint counts \(8,12\) \\
H2 & pure-frontier exclusion &
\((h_q\ge0\ \forall q)\wedge rawH\ge1171\) is \(\mathrm{UNSAT}\) \\
H3 & pure-frontier witnesses & \(\RP\) and \(\RC\) have
\(rawH=1170\), \(N_-=0\) \\
H4 & signed-cancellation spike & exactly twelve points above the pure
frontier, all with \((rawH,N_-)=(1217,3)\) \\
\end{longtable}

\begin{theorem}[Global frontier of \(\mathcal H\)]
\label{thm:7.2}
On
\(\Omega'=S^{21}\):

\begin{enumerate}
\def\labelenumi{\arabic{enumi}.}
\item
  The unrestricted range is \[
  \boxed{
  \min H_{tot}=-2268,
  \qquad
  \max H_{tot}=7302.
  }
  \] Equivalently, \(rawH\in[-378,1217]\). The endpoint multiplicities
  are \(8\) and \(12\).
\item
  The pure frontier is \[
  \boxed{
  \max\{H_{tot}:N_-=0\}=7020.
  }
  \]
\item
  The region above the pure frontier is exactly the unrestricted maximum
  layer: \[
  \boxed{
  \{H_{tot}>7020\}=\{H_{tot}=7302\}.
  }
  \] This layer has \(12\) points, all with \[
  rawH=1217,
  \qquad H_{tot}=7302,
  \qquad N_-=3,
  \] \[
  H_{pos}=7308,
  \qquad H_{\mathrm{neg\_abs}}=6,
  \qquad h_{\mathrm{loc,min}}=-2,
  \qquad h_{\mathrm{loc,max}}=18.
  \]
\end{enumerate}
\end{theorem}

\begin{proof}
The analytic reduction is Theorem~\ref{thm:7.1}. H1 covers all
\(3^{21}\) normal-form words and gives the unrestricted range. For the
pure frontier, \(N_-=0\) is the system \(h_q\ge0\) for all \(q\in S^5\),
and \(H_{tot}>7020\) is equivalent to \(rawH\ge1171\). H2 proves \[
(h_q\ge0\ \forall q\in S^5)\wedge rawH\ge1171
\] \(\mathrm{UNSAT}\) on all \(18540\) depth-9 live nodes: \[
18540/18540\ \mathrm{UNSAT},
\qquad \mathrm{SAT}=0,
\qquad \mathrm{UNKNOWN}=0.
\] H3 supplies pure witnesses at \(rawH=1170\), and H4 classifies the
full layer above the pure threshold as the twelve impure maximizers.
\end{proof}

\subsection{Pure frontier and signed-cancellation layer}\label{h-pure-frontier-spike}\label{rp-and-rc-on-the-pure-frontier}

\begin{corollary}
\label{cor:7.3}
\(\RP\) and \(\RC\) satisfy \[
H_{tot}=7020,
\qquad rawH=1170,
\qquad N_-=0.
\] Thus both lie on the global pure \(\mathcal H\)-frontier.
\end{corollary}

\begin{proof}
Direct evaluation of \[
\RP=111111111222222222000,
\qquad
\RC=222222222111111111000
\] gives the displayed values, and Theorem~\ref{thm:7.2} proves that no pure point
has larger \(H_{tot}\).
\end{proof}


\clearpage
\part{Intrinsic finite selection}

Part I gives the finite atlas data for the ambient landscape.  Part II
defines the selector and proves that its unique survivor is \(\RP\).  The
first two predicates are compression predicates, the third fixes the compressed
diagonal convention, and the fourth is the directed-edge geometry predicate.

\section{Selector predicates}\label{selector-predicates}

For a normal-form point of \(\Omega'\), write \(g\) for its cross-row output
column rule and \(d=(d_0,d_1,d_2)\) for its same-point diagonal.  The selector is
the conjunction
\[
\boxed{
\mathcal H_{acc}(g)=0,
\qquad
H_{diag}(d)=0,
\qquad
(r,r)^2=(r,r)\ (r\in S),
\qquad
\operatorname{AbsTrans}=\{C,J\}.
}
\]
The first predicate is hidden-column access zero.  The second is diagonal
Shannon entropy zero.  The third fixes the row-anchor representative after
entropy compression; equivalently, it anchors idempotence on the diagonal
orbit \(D_0\).  The fourth chooses the intrinsic head-row continuation/reversal
geometry on the six directed edges.

The first two predicates are information-theoretic: \(\mathcal H_{acc}\) is
conditional mutual information and \(H_{diag}\) is Shannon entropy
\cite{Shannon1948,CoverThomas2006}.  The Landauer--Bennett reading is the
implementation-independent economy reading of avoidable hidden column access
and discarded diagonal information
\cite{Landauer1961,Bennett1973,Bennett1982}.

\section{Access-zero and diagonal compression}\label{access-zero-diagonal-compression}

\subsection{Column-blind cross-row rules}\label{column-blind-rules}\label{column-blind-cross-row-rules}

For \(r_1\ne r_2\), write the cross-row branch as
\[
(r_1,c_1)\cdot(r_2,c_2)=(r_1,g(r_1,c_1,r_2,c_2)).
\]
A cross-row rule is \emph{column-blind} if the output column depends only on
its ordered row pair:
\[
g(r_1,c_1,r_2,c_2)=h(r_1,r_2).
\]
Equivalently, after \((r_1,r_2)\) is fixed, the output column is constant over
the nine choices of \((c_1,c_2)\).

By \(\sigma\)-equivariance, a column-blind rule is determined by two anchors
\[
a=h(0,1),
\qquad
b=h(0,2),
\]
and the remaining ordered row pairs satisfy
\[
h(r,r+1)=r+a,
\qquad
h(r,r+2)=r+b.
\]
There are \(3^2=9\) column-blind cross-row rules; denote them by
\(h_{a,b}\).

\subsection{Hidden-column access}\label{hidden-column-access-selector}\label{hidden-column-access}\label{hidden-column-access-as-a-landauer-facing-observable}

Let
\[
\Sigma_{cross}=\{(r_1,c_1,r_2,c_2)\in S^4:r_1\ne r_2\}
\]
with the uniform distribution, and let
\[
C_{out}=g(R_1,C_1,R_2,C_2)
\]
be the output-column random variable.  Define
\[
\boxed{
\mathcal H_{acc}(g)=I(C_{out};C_1,C_2\mid R_1,R_2).
}
\]
Logarithms are base \(2\).

\begin{theorem}[Access-zero theorem]
\label{thm:3.1}
For a cross-row rule \(g\),
\[
\mathcal H_{acc}(g)=0
\quad\Longleftrightarrow\quad
 g(r_1,c_1,r_2,c_2)=h(r_1,r_2).
\]
Thus the access-zero locus consists of the nine column-blind rules
\(h_{a,b}\).
\end{theorem}

\begin{proof}
Conditional mutual information vanishes exactly when
\[
C_{out}\perp (C_1,C_2)\mid (R_1,R_2).
\]
Since \(C_{out}\) is a deterministic function of
\((R_1,C_1,R_2,C_2)\), this conditional independence is equivalent to
constancy in \((C_1,C_2)\) for every fixed ordered row pair.  This is exactly
column-blindness.  The ordered row-pair representatives \((0,1)\) and
\((0,2)\) then supply the two anchors \((a,b)\in S^2\).
\end{proof}

\subsection{Diagonal entropy}\label{diagonal-entropy-selector}\label{diagonal-compression}

For a diagonal triple \(d=(d_0,d_1,d_2)\in S^3\), let
\[
p_s=\frac13\#\{i:d_i=s\}.
\]
The diagonal Shannon observable is
\[
\boxed{
H_{diag}(d)=-\sum_{s\in S}p_s\log_2 p_s.
}
\]
The possible levels are
\[
\begin{array}{c|c|c}
\text{diagonal type} & \text{multiplicities} & H_{diag} \\
\hline
\text{constant} & (3,0,0) & 0 \\
\text{two-valued} & (2,1,0) & \log_2 3-\frac23 \\
\text{permutation} & (1,1,1) & \log_2 3
\end{array}
\]

\begin{theorem}[Diagonal compression]
\label{thm:9.1}
For diagonal triples,
\[
H_{diag}(d)=0
\quad\Longleftrightarrow\quad
 d_0=d_1=d_2.
\]
Thus diagonal compression reduces the diagonal sector from \(27\) choices to
\[
000,
\qquad 111,
\qquad 222.
\]
Together with access-zero, the selector domain becomes \(9\times3\).
\end{theorem}

\begin{proof}
Entropy of a finite distribution vanishes exactly when the distribution is
concentrated on one symbol.  For the empirical distribution of
\((d_0,d_1,d_2)\), this is precisely constancy of the diagonal triple.
\end{proof}

\subsection{Row-anchor idempotence}\label{row-anchor-idempotence}

The row anchors are the points of the diagonal orbit \(D_0\).

\begin{proposition}[Row-anchor idempotence]
\label{prop:10.2}
After diagonal compression, the condition
\[
(r,r)\cdot(r,r)=(r,r)
\qquad(r\in S)
\]
selects the unique constant diagonal
\[
d=000.
\]
For this diagonal,
\[
(r,r+i)^2=(r,r)
\qquad(r,i\in S),
\]
so the idempotent points are exactly the diagonal orbit \(D_0\).
\end{proposition}

\begin{proof}
On the compressed diagonal sector \(d=sss\), the equation
\((r,r)^2=(r,r)\) gives \(s=0\).  Hence the unique compressed diagonal
compatible with row-anchor idempotence is \(000\).  The displayed square formula
follows by substituting \(d=000\) into the same-point branch.
\end{proof}

The first three selector predicates therefore reduce the normal-form
coordinates to
\[
h_{a,b},
\qquad
(a,b)\in S^2,
\qquad
 d=000.
\]

\section{Directed-edge geometry}\label{directed-edge-geometry}

\subsection[The maps C and J]{The maps \(C\) and \(J\)}\label{maps-c-j}

Let
\[
M^\times=\{(r,c)\in M:r\ne c\}
\]
be the six-point directed-edge set of the triangle on vertex set \(S\).  In the
orbit notation above,
\[
M^\times=D_1\sqcup D_2.
\]
The point \((r,c)\) is the directed edge with tail \(r\) and head \(c\).  Define
\[
C(r,c)=(c,\overline{rc}),
\qquad
J(r,c)=(c,r).
\]
The map \(C\) advances the edge through the third vertex, while \(J\) reverses
it.

\begin{theorem}[Directed-edge dihedral relations]
\label{thm:11.1}
On \(M^\times\),
\[
C^3=1,
\qquad
J^2=1,
\qquad
JCJ=C^{-1}.
\]
Thus \(C\) and \(J\) generate the dihedral geometry of the directed triangle.
\end{theorem}

\begin{proof}
Applying \(C\) three times gives
\[
(r,c)\mapsto(c,\overline{rc})\mapsto(\overline{rc},r)\mapsto(r,c).
\]
Applying \(J\) twice reverses the edge twice.  Finally,
\[
JCJ(r,c)=JC(c,r)=J(r,\overline{rc})=(\overline{rc},r)=C^{-1}(r,c).
\]
\end{proof}

For every relabelling \(\rho\in\operatorname{Sym}(S)\), acting by
\(\rho_\times(r,c)=(\rho(r),\rho(c))\), one has
\[
\rho_\times C=C\rho_\times,
\qquad
\rho_\times J=J\rho_\times.
\]
The pair \(\{C,J\}\) is the intrinsic head-row pair:
\[
\{C(r,c),J(r,c)\}
=
\{y\in M^\times: \operatorname{tail}(y)=c\}.
\]

\subsection{Absorption transforms}\label{absorption-transforms}\label{absorption-and-absorption-transitions}

For \(x\in M^\times\), define
\[
\operatorname{Abs}(x)=\{y\in M^\times\setminus\{x\}:x\cdot y=x\}.
\]
On the anchored selector domain \(h_{a,b}\), \(d=000\), absorption of
\(x=(r,c)\) is governed by the cross-row equation
\[
h(r,r_y)=c.
\]
The pure \(C/J\) predicate is the finite condition
\[
\operatorname{Abs}(x)=\{C(x),J(x)\}
\qquad(x\in M^\times).
\]
Equivalently, the anchored rule has two absorbed-neighbour maps, and they form
\[
\operatorname{AbsTrans}=\{C,J\}.
\]

\subsection[Pure C/J]{Pure \(C/J\)}\label{pure-cj-selector}\label{pure-cj-and-the-selection-theorem}

\begin{theorem}[Pure \(C/J\) selector on the anchored domain]
\label{thm:12.2}
For a column-blind rule \(h_{a,b}\) with \(d=000\),
\[
\operatorname{AbsTrans}=\{C,J\}
\quad\Longleftrightarrow\quad
(a,b)=(1,2).
\]
The corresponding cross-row branch is the right-row projection branch.
\end{theorem}

\begin{proof}
Let \(x=(r,c)\in M^\times\), and write \(\delta=c-r\in\{1,2\}\).  For the two
cross rows of a potential absorber, the absorption equation becomes
\[
r_y=r+1\Rightarrow a=\delta,
\qquad
r_y=r+2\Rightarrow b=\delta.
\]
The pure \(C/J\) predicate requires the absorbed neighbours of \((r,c)\) to be
exactly
\[
J(r,c)=(c,r),
\qquad
C(r,c)=(c,\overline{rc}),
\]
so both absorbed neighbours must lie in row \(c\).

For \(\delta=1\), this forces absorption in row \(r+1=c\), hence \(a=1\), and
excludes absorption in row \(r+2\), hence \(b\ne1\).  For \(\delta=2\), it
forces absorption in row \(r+2=c\), hence \(b=2\), and excludes absorption in
row \(r+1\), hence \(a\ne2\).  Therefore \((a,b)=(1,2)\).

Conversely, if \((a,b)=(1,2)\), then for every \(x=(r,c)\) the only absorbing
cross row is row \(c\), and the two off-diagonal points in that row are
\((c,r)\) and \((c,\overline{rc})\), namely \(J(x)\) and \(C(x)\).  Thus
\(\operatorname{AbsTrans}=\{C,J\}\).
\end{proof}

\section{Uniqueness theorem}\label{uniqueness-theorem}\label{compact-uniqueness-core}

\begin{theorem}[Uniqueness theorem]
\label{thm:12.3}
The four conditions
\[
\mathcal H_{acc}=0,
\qquad
H_{diag}=0,
\qquad
(r,r)^2=(r,r)\ (r\in S),
\qquad
\operatorname{AbsTrans}=\{C,J\}
\]
select \(\RP\) uniquely inside \(\Omega'\).
\end{theorem}

\begin{proof}
Access-zero gives a column-blind rule \(h_{a,b}\) by
Theorem~\ref{thm:3.1}.  Diagonal entropy zero gives a constant diagonal by
Theorem~\ref{thm:9.1}, and row-anchor idempotence fixes it to \(d=000\) by
Proposition~\ref{prop:10.2}.  On this anchored domain, the pure \(C/J\)
selector is equivalent to \((a,b)=(1,2)\) by Theorem~\ref{thm:12.2}.  This is
exactly the right-row projection multiplication, with
\[
(r_1,c_1)\cdot(r_2,c_2)=(r_1,r_2)
\qquad(r_1\ne r_2),
\]
and same-point diagonal \(d=000\).  Hence the survivor is unique and is
\(\RP\).
\end{proof}

\clearpage

\part{Linearization of the selected magma}

Part III fixes the selected multiplication \(\RP\) and records the structures
that appear after this choice.  The first layer is algebraic: row ownership,
left translations, composition envelope, row-fiber algebra, endomorphisms,
intertwiners, and binary term functions.  The second layer is carrier-level:
row-contraction on \(k[M]\cong\operatorname{Mat}_3(k)\), the
hidden-continuation operator bridge, compact companions, and the signed
symplectic lift of the pure directed-edge geometry.

\section{The selected operation}\label{selected-operation}\label{the-selected-rp-operation}

\subsection{Multiplication and ground-field conventions}\label{selected-multiplication-ground-fields}\label{rp-multiplication-and-ground-fields}

The selected multiplication is
\[
(r_1,c_1)\cdot(r_2,c_2)=
\begin{cases}
(r_1,r_2), & r_1\ne r_2,\\[2mm]
(r_1,\overline{c_1c_2}), & r_1=r_2,\ c_1\ne c_2,\\[2mm]
(r_1,r_1), & r_1=r_2,\ c_1=c_2.
\end{cases}
\]
The first case is the right-row projection branch; the second is the fixed
same-row Steiner branch; the third is diagonal anchoring.

The finite magma statements are set-theoretic.  When the carrier is
linearized, let \(k\) be a ground field and let \(k[M]\) be the vector space
with basis \(e_x\), \(x\in M\).  Rank and image statements for individual
finite maps are valid over any ground field.  Linear span statements are over
the chosen \(k\).  Spectral, Jordan, and fiber-algebra statements involving
\(t\pm1\) or idempotents with coefficients \(1/2\) are stated under
\[
\operatorname{char}k\ne2.
\]
The carrier-level compact-companion and signed symplectic bridge statements
are over \(\mathbb R\) unless stated otherwise.
\subsection{Row-owner law}\label{row-owner-law}\label{row-owner-principle}

For a nonempty term \(t(x_1,\ldots,x_n)\), define its \textbf{owner}
\(\operatorname{ow}(t)\) to be the leftmost variable appearing in the term
tree.

\begin{lemma}[Row-owner lemma]
\label{lem:13.1}
For every \(\RP\) term \(t(x_1,\ldots,x_n)\),
\[
\operatorname{row}\bigl(t(x_1,\ldots,x_n)\bigr)
=
\operatorname{row}\bigl(\operatorname{ow}(t)\bigr).
\]
\end{lemma}

\begin{proof}
The assertion is immediate for a variable.  If \(t=u\cdot v\), then the
left-row output axiom gives
\(\operatorname{row}(t)=\operatorname{row}(u\cdot v)=\operatorname{row}(u)\).
The induction hypothesis identifies this row with the row of the leftmost
variable of \(u\), which is also the leftmost variable of \(u\cdot v\).
\end{proof}

The row-owner lemma is the common source of the row-target product law for
left-translation compositions and the owner coordinate in the binary term
algebra.

\section{Left translations and operator envelope}\label{left-translations-operator-envelope}\label{left-regular-operators}

For \(x\in M\), write \(L_x(y)=x\cdot y\).

\subsection{Rank and image}\label{rank-and-image}\label{row-fiber-normal-form}

\begin{proposition}
\label{prop:14.1}
For \(x=(r,c)\),
\[
L_{(r,c)}(s,d)=
\begin{cases}
(r,s), & s\ne r,\\[1mm]
(r,r), & s=r,\ d=c,\\[1mm]
(r,\overline{cd}), & s=r,\ d\ne c.
\end{cases}
\]
Consequently \(L_{(r,c)}(M)=F_r\).
\end{proposition}

\begin{proof}
The three cases are exactly the cross-row projection, diagonal anchoring, and
same-row Steiner branches of \(\RP\).  They produce all three points of
\(F_r\).
\end{proof}

\begin{theorem}
\label{thm:14.2}
For every \(x\in M\),
\[
\operatorname{rank}L_x=3,
\qquad
\operatorname{Im}L_x=F_{\operatorname{row}(x)}.
\]
Moreover,
\[
\dim\operatorname{span}\{L_x:x\in M\}=9.
\]
\end{theorem}

\begin{proof}
Proposition~\ref{prop:14.1} shows that each \(L_x\) maps onto the three-point
row fiber \(F_{\operatorname{row}(x)}\), hence has rank three after
linearization.  Operators with different target rows have disjoint codomain
support.  For a fixed row, the three restrictions to the row fiber are
linearly independent; in normalized row \(0\) they contain the matrices
\[
\begin{pmatrix}1&0&0\\0&0&1\\0&1&0\end{pmatrix},
\qquad
\begin{pmatrix}0&1&1\\0&0&0\\1&0&0\end{pmatrix},
\qquad
\begin{pmatrix}0&1&1\\1&0&0\\0&0&0\end{pmatrix}.
\]
Thus each row contributes dimension \(3\), and the total dimension is \(9\).
The \texttt{L3/} operator rank and spectrum audit tables record the nine-row
calculation.
\end{proof}

Assume \(\operatorname{char}k\ne2\).

\begin{theorem}
\label{thm:14.3}
If \(x=(r,r)\), then
\[
\mu_{L_x}(t)=t(t-1)(t+1)
\]
and \(L_x\) is semisimple.  If \(x=(r,c)\) with \(r\ne c\), then
\[
\mu_{L_x}(t)=t^2(t-1)(t+1),
\]
with zero-primary Jordan structure \(J_2(0)\oplus J_1(0)^{\oplus5}\).
\end{theorem}

\begin{proof}
Every \(L_{(r,c)}\) maps \(k[M]\) onto \(kF_r\), so the nonzero spectral
part is the restriction to \(kF_r\).  For \(c=r\), this restriction fixes
\((r,r)\) and swaps the two off-diagonal basis vectors.  For \(c\ne r\),
the restriction has eigenvalues \(0,1,-1\), and its one-dimensional kernel
lies inside the image, producing exactly one length-two zero block.  The
operator-spectrum table audits the calculation for all nine translations.
\end{proof}

\subsection{Positive composition envelope}\label{positive-composition-envelope}\label{composition-envelope-and-directed-edge-representation}

Let \(\operatorname{Mon}^+\langle L_x\rangle\) be the positive composition
envelope generated by the nine left translations.  The superscript \(+\) means
nonempty words in the generators; the identity map is not included.  It is a
semigroup under composition, and the unital monoid obtained by adjoining the
identity has one additional element.  For \(r\in S\), let \(\mathcal W_r\) be
the component of maps whose image lies in \(F_r\).

For fixed \(r\), define
\[
\tau_{r,c}(d)=
\begin{cases}
r, & d=c,\\
\overline{cd}, & d\ne c,
\end{cases}
\qquad
\mathcal P_r=\langle \tau_{r,c}:c\in S\rangle.
\]
In normalized row \(0\),
\[
\mathcal P_0=\{012,021,200,100,011,022\}.
\]
For \(\psi\in\mathcal P_r\) and \(c\in S\), set
\[
N_{r,\psi,c}(s,d)=
\begin{cases}
(r,\psi(\tau_{r,c}(d))), & s=r,\\
(r,\psi(s)), & s\ne r,
\end{cases}
\]
and let \(K_{r,t}\) be the constant map to \((r,t)\).

\begin{proposition}[Fixed-row composition normal form]
\label{prop:15.1}
For every \(r\in S\),
\[
\mathcal W_r=
\{N_{r,\psi,c}:\psi\in\mathcal P_r,\ c\in S\}
\cup
\{K_{r,t}:t\in S\}.
\]
There are \(14\) nonconstant maps of the first type and \(3\) constants;
hence \(|\mathcal W_r|=17\).
\end{proposition}

\begin{proof}
A positive composition is owned by its leftmost generator.  If all later
generators have the same target row \(r\), the word has the displayed
\(N_{r,\psi,c}\) form.  If some later generator has target row different from
\(r\), the leftmost row-\(r\) generator sees a cross-row input and the
composition becomes constant.  The normalized finite count \(18-4=14\) for
the nonconstant forms is audited in the composition envelope table.
\end{proof}

\begin{theorem}
\label{thm:15.2}
The positive composition envelope generated by the left translations has size
\[
|\operatorname{Mon}^+\langle L_x\rangle|=51.
\]
The identity map is not contained in it; adjoining the identity gives the
unital generated monoid of size \(52\).  More precisely,
\[
\operatorname{Mon}^+\langle L_x\rangle
=
\mathcal W_0\sqcup\mathcal W_1\sqcup\mathcal W_2,
\qquad |\mathcal W_r|=17.
\]
Its linear envelope decomposes as
\[
\mathcal E=W_0\oplus W_1\oplus W_2,
\qquad
\dim W_r=9,
\qquad
\dim\mathcal E=27,
\]
and satisfies \(W_rW_s\subseteq W_r\).
\end{theorem}

\begin{proof}
Proposition~\ref{prop:15.1} gives \(|\mathcal W_r|=17\), and the three
target-row components are disjoint.  Every positive word has image contained in
a single row fiber, so the identity map on \(M\) is not a positive word.  The
row-owner law gives \(W_rW_s\subseteq W_r\).  The linear dimension claim is the
row-wise basis-reduction audit recorded in the \texttt{L3/} composition-envelope
and product-law tables.
\end{proof}

\subsection{Row-fiber algebra}\label{row-fiber-algebra}\label{the-five-dimensional-row-fiber-algebra}

Assume \(\operatorname{char}k\ne2\).  For each row \(r\), define
\[
A_{5,r}=\langle L_{(r,c)}|_{kF_r}:c\in S\rangle
\subseteq\operatorname{End}(kF_r).
\]

\begin{proposition}
\label{prop:15.3}
The algebra \(A_{5,r}\) has dimension \(5\).  Its center is \(kI\), its
radical is two-dimensional and square-zero, and
\[
A_{5,r}/\operatorname{Rad}(A_{5,r})\cong k^3.
\]
For \(r=0\), one basis is
\[
L_{(0,0)}|_{kF_0},\quad L_{(0,1)}|_{kF_0},\quad L_{(0,2)}|_{kF_0},
\quad I,\quad (L_{(0,1)}|_{kF_0})^2.
\]
\end{proposition}

\begin{proof}
By \(\sigma\)-equivariance it is enough to work in row \(0\).  The
five-dimensional multiplication table, center, radical, quotient, and
operator-boundary checks are recorded in the auxiliary fiber-algebra
certificate files listed in \texttt{SUPPORT\_INDEX.txt}.
\end{proof}

\begin{theorem}
\label{thm:15.4}
On the six directed edges,
\[
\mathbb R[M^\times]\cong\mathbf1\oplus\mathrm{sign}\oplus2\,\mathrm{std}.
\]
For \(\RP\), the absorption matrix is
\[
A_{abs}=(P_C+P_J)^\top
\]
and has spectrum
\[
\{2,0,0,0,-1,-1\}.
\]
\end{theorem}

\begin{proof}
The six directed edges carry the permutation representation of the directed
triangle.  The pure absorption transitions of \(\RP\) are exactly \(C\) and
\(J\), so the absorption matrix is \((P_C+P_J)^\top\).  The transition and
spectrum tables audit the finite calculation.
\end{proof}

\section{Endomorphisms and intertwiners}\label{endomorphisms-intertwiners}\label{endomorphisms-intertwiners-and-row-contraction}

For a map \(\rho:S\to S\), write
\(\varphi_\rho(r,c)=(\rho(r),\rho(c))\).  For \(t\in S\), write
\(e_t(r,c)=(t,t)\).

For a row \(r\), set
\[
c\star_r d=
\begin{cases}
r, & c=d,\\
\overline{cd}, & c\ne d.
\end{cases}
\]

\begin{lemma}[Row-fiber homomorphisms]
\label{lem:16.1}
If \(\phi:S\to S\) satisfies \(\phi(r)=t\) and
\[
\phi(c\star_r d)=\phi(c)\star_t\phi(d)
\qquad(c,d\in S),
\]
then \(\phi\) is either the constant map to \(t\), or a permutation of
\(S\) sending \(r\) to \(t\).
\end{lemma}

\begin{proof}
After translating to \(r=t=0\), write \(\phi(1)=u\), \(\phi(2)=v\).  The
relations \(0\star_0 1=2\) and \(0\star_0 2=1\) force
\[
(u,v)=(0,0),\quad (1,2),\quad (2,1),
\]
which are respectively the constant solution and the two permutations fixing
\(0\).
\end{proof}

\begin{theorem}
\label{thm:16.2}
The endomorphism monoid of the selected \(\RP\) magma is
\[
\operatorname{End}(M,\cdot)
=
\{\varphi_\rho:\rho\in\operatorname{Sym}(S)\}\cup\{e_0,e_1,e_2\}.
\]
\end{theorem}

\begin{proof}
The idempotents of \(\RP\) are exactly the diagonal points \((r,r)\).  Thus
an endomorphism sends row fibers to row fibers over a map \(\rho:S\to S\),
and its restriction to each row fiber is governed by Lemma~\ref{lem:16.1}.
The cross-row rule forces all nonconstant row restrictions to be the same
permutation \(\rho\).  If \(\rho\) fails to be injective, the cross-row
equations force all affected row restrictions to collapse, and then the third
row collapses as well; the map is one of the \(e_t\).  Conversely,
simultaneous permutations and diagonal collapses preserve the three branches
of the \(\RP\) multiplication.
\end{proof}

For an endomorphism \(\varphi\), let \(\operatorname{Hom}_\varphi\) be the
space of intertwiners satisfying
\[
T L_x=L_{\varphi(x)}T
\qquad(x\in M).
\]
Let \(V=k[M]\), and let
\[
\varepsilon:V\to k,
\qquad
\varepsilon\!\left(\sum_{m\in M}a_m e_m\right)=\sum_{m\in M}a_m
\]
be the total augmentation functional.  For \(r\in S\), set
\[
d_r=e_{(r,r)},
\qquad
o_r=\sum_{c\ne r}e_{(r,c)},
\qquad
U_r=\operatorname{span}\{d_r,o_r\}.
\]
Here \(u\varepsilon\) denotes the rank-one operator
\(v\mapsto \varepsilon(v)u\).

\begin{proposition}[Constant intertwiners]
\label{prop:constant-intertwiners}
For each diagonal collapse \(e_r(x)=(r,r)\),
\[
\operatorname{Hom}_{e_r}
=
\{u\varepsilon:u\in U_r\}
=
\left\{
\left(\alpha e_{(r,r)}+\beta\sum_{c\ne r}e_{(r,c)}\right)\varepsilon:
\alpha,\beta\in k
\right\}.
\]
In particular, \(\dim\operatorname{Hom}_{e_r}=2\).
\end{proposition}

\begin{proof}
Each \(L_x\) is induced by a map \(M\to M\), hence every column of its matrix
has sum one and \(\varepsilon L_x=\varepsilon\).  If
\(u\in\operatorname{Fix}(L_{(r,r)})\), then
\[
(u\varepsilon)L_x=u\varepsilon
\qquad\text{and}\qquad
L_{(r,r)}(u\varepsilon)=(L_{(r,r)}u)\varepsilon=u\varepsilon,
\]
so \(u\varepsilon\in\operatorname{Hom}_{e_r}\).

Conversely, if \(T\in\operatorname{Hom}_{e_r}\), then
\[
T L_x=L_{(r,r)}T=T L_y
\qquad(x,y\in M),
\]
so \(T\) kills
\[
\mathcal D=\operatorname{span}\{(L_x-L_y)v:x,y\in M,\ v\in V\}.
\]
Because \(\varepsilon L_x=\varepsilon\), one has
\(\mathcal D\subseteq\ker\varepsilon\).  The reverse inclusion is finite:
the graph generated by relations \(x\cdot z\sim y\cdot z\), with common
source \(z\), is connected.  The Front-C verifier records an explicit
eight-edge spanning tree, so the differences in \(\mathcal D\) span the
augmentation-zero hyperplane.  Thus \(\mathcal D=\ker\varepsilon\), and
\(T\) factors uniquely as \(T=u\varepsilon\).

Substitution into the intertwining equation gives \(L_{(r,r)}u=u\).  Finally,
\(L_{(r,r)}\) has image in the row fiber
\(F_r=\operatorname{span}\{e_{(r,c)}:c\in S\}\); on this fiber it fixes
\(e_{(r,r)}\) and swaps the two off-diagonal basis vectors.  Therefore
\[
\operatorname{Fix}(L_{(r,r)})
=
\operatorname{span}\left\{e_{(r,r)},\sum_{c\ne r}e_{(r,c)}\right\}=U_r,
\]
which proves the formula.
\end{proof}

\begin{proposition}[Intertwiner classification]
\label{prop:intertwiner-classification}
Let \(P_\rho\) be the permutation operator on \(V\) induced by
\(\varphi_\rho\).  For \(\rho\in\operatorname{Sym}(S)\),
\[
\operatorname{Hom}_{\varphi_\rho}=kP_\rho,
\qquad
\dim\operatorname{Hom}_{\varphi_\rho}=1.
\]
For the diagonal collapse \(e_r\),
\[
\operatorname{Hom}_{e_r}=U_r\varepsilon,
\qquad
\dim\operatorname{Hom}_{e_r}=2.
\]
Thus the total audited intertwiner dimension is
\[
6\cdot1+3\cdot2=12.
\]
\end{proposition}

\begin{proof}
For an automorphism \(\varphi_\rho\), the operator \(P_\rho\) satisfies
\[
P_\rho L_x=L_{\varphi_\rho(x)}P_\rho
\qquad(x\in M),
\]
so \(kP_\rho\subseteq\operatorname{Hom}_{\varphi_\rho}\).  The Front-C verifier
audits the defining linear constraints and gives
\(\dim\operatorname{Hom}_{\varphi_\rho}=1\) for the six automorphism spaces;
hence equality holds.  The formula for \(\operatorname{Hom}_{e_r}\) and its
dimension are Proposition~\ref{prop:constant-intertwiners}.  Summing over the
six automorphisms and three diagonal collapses gives the displayed total.
\end{proof}

\section{Binary term functions}\label{binary-term-functions}\label{binary-term-algebra}

A binary term function of \(\RP\) is a function \(M^2\to M\) realized by a
term in variables \(x,y\).  The set of all such functions is denoted
\(T_2(\RP)\).

By the row-owner lemma, every binary term has an owner variable
\(\operatorname{ow}(f)\in\{x,y\}\).  Order the inputs as
\((o,n)=(\operatorname{owner},\operatorname{nonowner})\).  On the same-row
sector, normalize the common row to \(0\) and define
\[
s_f(i,j)=\operatorname{col}\bigl(f(o=(0,i),n=(0,j))\bigr).
\]
On the cross-row sector, normalize the owner row to \(0\) and the nonowner
row to \(1\); the output column is independent of the nonowner column, so
\[
u_f(i)=\operatorname{col}\bigl(f(o=(0,i),n=(1,j))\bigr).
\]
The opposite cross-row direction is determined by \(u_f^{-}(i)=-u_f(-i)\).

\begin{proposition}[Binary sector reduction]
\label{prop:18.1}
Every binary term function has a normal signature
\[
\Theta(f)=\bigl(\operatorname{ow}(f),s_f,u_f\bigr).
\]
The same-row signature \(s_f:S^2\to S\) and cross unary signature
\(u_f:S\to S\) determine the term function on all row sectors.
\end{proposition}

\begin{proof}
The owner coordinate is Lemma~\ref{lem:13.1}.  Same-row inputs remain in the
common input row.  Cross-row inputs are governed by the owner row and a unary
owner-column signature; the opposite direction is forced by the
\(\RP\)-compatible relabelling exchanging the two nonowner directions.
\end{proof}

The realized same-row signatures form a \(21\)-element set \(\mathcal S_{21}\).
The realized cross unary signatures form a \(15\)-element set
\(\mathcal U_{15}\).  The candidate normal forms are
\[
\{x,y\}\times\mathcal S_{21}\times\mathcal U_{15}.
\]

\begin{lemma}[Binary term certificate protocol]
\label{lem:18.2}
The binary term theorem is supported by four finite obligations: T2-SIG exposes
\(\mathcal S_{21}\) and \(\mathcal U_{15}\); T2-REAL realizes every triple;
T2-TRANS checks closure of the two factor transition systems; T2-MULT
reconstructs the full \(630\times630\) multiplication table from the owner and
factor transitions.
\end{lemma}

\begin{theorem}
\label{thm:18.3}
The binary term-function algebra of \(\RP\) is finite and satisfies
\[
T_2(\RP)
\cong
\{x,y\}\times\mathcal S_{21}\times\mathcal U_{15}.
\]
Consequently
\[
|T_2(\RP)|=2\cdot21\cdot15=630.
\]
\end{theorem}

\begin{proof}
Proposition~\ref{prop:18.1} gives injectivity of the normal signature.  The
realization certificate supplies one realizing term for every triple, and the
finite transition tables reconstruct the product law.  Hence the cardinality
is \(2\cdot21\cdot15\).
\end{proof}

This is the binary layer of the term algebra.  The all-arity clone is left for
the open-problem list.

\section{Matrix carrier and row-contraction}\label{matrix-carrier-row-contraction}\label{cross-row-row-contraction}

Under the carrier identification
\[
k[M]\cong\operatorname{Mat}_3(k),
\qquad e_{(r,c)}\leftrightarrow E_{rc},
\]
the finite \(\RP\) multiplication and ordinary matrix multiplication become
two bilinear structures on the same nine-dimensional vector-space carrier.
The row-contraction bridge records the cross-row branch of \(\RP\) on matrix
units.

\begin{proposition}
\label{prop:16.3}
Under \(k[M]\cong\operatorname{Mat}_3(k)\), the row-contraction extension of
the cross-row \(\RP\) rule is
\[
\mu_{cross}(A,B)=AJ_3B^\top=(A\mathbf1)(B\mathbf1)^\top.
\]
On matrix units with cross rows,
\[
E_{ij}J_3E_{kl}^\top=E_{ik}.
\]
\end{proposition}

\begin{proof}
Since \(J_3=\mathbf1\mathbf1^\top\),
\[
E_{ij}J_3E_{kl}^\top
=(e_ie_j^\top)(\mathbf1\mathbf1^\top)(e_le_k^\top)
=E_{ik}.
\]
For \(i\ne k\), this is exactly the cross-row \(\RP\) rule
\((i,j)\cdot(k,l)=(i,k)\).  The \texttt{L3/} matrix-bridge audit tables
record the basis-level bridge.
\end{proof}

Thus the matrix bridge expresses the same column-blindness selected by
\(\mathcal H_{acc}=0\): cross-row multiplication factors through row sums.
The same-row branch remains the Steiner/diagonal row-fiber rule.

\section{Hidden-continuation operator bridge}\label{hidden-continuation-operator-bridge}\label{hidden-continuation-operator-identity}

The local hidden-continuation identity of Section~\ref{hidden-continuation-frontier}
has an operator interpretation in the left-regular envelope.  For a normalized
local index \(q\), \(D_C\) compares the two stepwise continuation signatures,
while \(D_L\) compares the two collapsed left translations.

\begin{theorem}[Hidden-continuation operator identity]
\label{thm:17.1}
Let \(g\in\Omega'\) be a completed normal-form magma and let
\(q=(b,t,a,e,f)\in S^5\).  Define
\[
xy,
\qquad
 yz,
\qquad
 z_R,
\qquad
 ep_L,
\qquad
 ep_R
\]
from \(g\) as in Theorem~\ref{thm:7.1}.  Then
\[
\frac{h_q(g)}2
=
D_C(t,xy,f;b,a,z_R)-D_L(ep_L,ep_R).
\]
For brevity write
\[
D_C(q;g)=D_C(t,xy,f;b,a,z_R),
\qquad
D_L(q;g)=D_L(ep_L,ep_R).
\]

After selection, take \(g=\RP\) and write the normalized triple as
\[
X=(0,a),\qquad Y=(b,e),\qquad Z=(t,f).
\]
Then
\[
XY=(0,xy),\qquad YZ=(b,z_R),
\]
and
\[
(XY)Z=(0,ep_L),\qquad X(YZ)=(0,ep_R).
\]
With the convention \(AB=A\circ B\), the continuation distance
\(D_C(q;\RP)\) is the Hamming-distance contrast between the normalized
signatures of
\[
L_{XY}L_Z
\quad\text{and}\quad
L_XL_{YZ},
\]
whereas \(D_L(q;\RP)\) is the Hamming-distance contrast between
\[
L_{(XY)Z}
\quad\text{and}\quad
L_{X(YZ)}.
\]
As maps, all four operators lie in \(\mathcal W_{\operatorname{row}(X)}\);
after linearization they lie in \(W_{\operatorname{row}(X)}\).

Define
\[
I_{tot}(g)=6\sum_{q\in S^5}D_C(q;g),
\qquad
B_{tot}(g)=6\sum_{q\in S^5}D_L(q;g).
\]
Thus \(H_{tot}(g)=I_{tot}(g)-B_{tot}(g)\).  For \(\RP\),
\[
I_{tot}=8904,
\qquad
B_{tot}=1884,
\qquad
H_{tot}=7020,
\qquad
N_-=0.
\]
The finalist \(\RC\) has the same pure hidden-continuation profile.
\end{theorem}

\begin{proof}
The scalar identity is exactly Theorem~\ref{thm:7.1} divided by \(2\):
\(D_C\) is the sum of the nine continuation disagreement indicators,
\(D_L\) is the sum of the nine collapsed left-translation disagreement
indicators, and \(h_q\) is twice their signed difference.

For \(g=\RP\), the displayed normalized triple identifies the two
continuation signatures with the signatures of \(L_{XY}L_Z\) and
\(L_XL_{YZ}\), under the convention \(AB=A\circ B\).  Likewise the collapsed
signatures are the left translations by the endpoints \((XY)Z\) and
\(X(YZ)\).  The row-target assertion follows from the composition-envelope law
\(W_rW_s\subseteq W_r\).  The \texttt{L3/} hidden-continuation bridge audit
records the \(\RP\)/\(\RC\) comparison and verifies the displayed values.  The
global frontier status is Theorem~\ref{thm:7.2}.
\end{proof}

\section{Carrier-level compact companions}\label{carrier-level-compact-companions}\label{carrier-level-companion-lie-package}

The carrier identification
\[
k[M]\cong\operatorname{Mat}_3(k),
\qquad e_{(r,c)}\leftrightarrow E_{rc},
\]
places the selected finite magma on the same nine-dimensional carrier as the
ordinary matrix algebra.  A \emph{carrier-level compact companion} is a compact
Lie group attached to a specified linear structure on this shared carrier.  The
ordinary matrix and Lie structures used here are standard
\cite{HornJohnson2012,FultonHarris1991,Hall2015}.

Over fields of characteristic different from \(3\), ordinary matrix
multiplication gives
\[
\mathfrak{gl}(3,k)=kI\oplus\mathfrak{sl}(3,k),
\qquad
\operatorname{Mat}_3(k)=kI\oplus\mathfrak{sl}(3,k)
\]
as vector spaces.  After complex compact form, this is the usual central
\(\mathfrak u(1)\) plus \(\mathfrak{su}(3)\) picture.

Over \(\mathbb R\), the compact three-dimensional branch
\[
\mathfrak{so}(3)=\{X\in\mathfrak{sl}(3,\mathbb R):X^\top=-X\}
\]
is stable under simultaneous permutation conjugation.  A split comparator is
\(\mathfrak{sl}(V_{std})\subset\mathfrak{sl}(3,\mathbb R)\), where
\[
V_{std}=\{v:v_0+v_1+v_2=0\}.
\]
Compactness selects the \(\mathfrak{so}(3)\) branch among these
three-dimensional \(\operatorname{Sym}(S)\)-stable carrier branches.

\begin{theorem}
\label{thm:19.1}
In the real/complex compact-form setting, the carrier identification and the
\(\RP\)-compatible \(\operatorname{Sym}(S)\)-action expose the carrier-level
compact companions
\[
SU(3),
\qquad
SU(2),
\qquad
U(1).
\]
Here \(SU(3)\) is attached to the ordinary \(\mathfrak{sl}(3)\) matrix
package, \(SU(2)\) is the simply connected compact companion of the selected
\(\mathfrak{so}(3)\) branch, and \(U(1)\) is attached to the central scalar
line.
\end{theorem}

\begin{proof}
The matrix-unit carrier gives \(\operatorname{Mat}_3\), hence the ordinary
\(\mathfrak{sl}(3)\) matrix package and its compact form \(SU(3)\).  The
\(\operatorname{Sym}(S)\)-stable compact branch \(\mathfrak{so}(3)\) gives
\(SO(3)\), whose simply connected compact cover is \(SU(2)\).  The central
scalar line gives the compact circle \(U(1)\).  The \texttt{L3/}
companion-structure bridge audit records these carrier-level compact types.
\end{proof}

\section{Signed symplectic lift}\label{signed-symplectic-lift}\label{signed-symplectic-bridge}

This bridge uses the finite pure \(C/J\) directed-edge geometry selected in
Part II.  The continuous statement is a signed projective symplectic
realization of that six-edge geometry and completes the carrier-level bridge
part of the passport \cite{McDuffSalamon2017}.

Let \(V=\mathbb R^6\) have ordered basis
\[
(e_{01},e_{12},e_{20},e_{10},e_{02},e_{21})
\]
and symplectic matrix
\[
\Omega=
\begin{pmatrix}0&I_3\\-I_3&0\end{pmatrix}.
\]
Set
\[
P=\begin{pmatrix}0&0&1\\1&0&0\\0&1&0\end{pmatrix},
\qquad
R=\begin{pmatrix}1&0&0\\0&0&1\\0&1&0\end{pmatrix},
\]
\[
\widehat C=\begin{pmatrix}P&0\\0&P\end{pmatrix},
\qquad
\widehat J_0=\begin{pmatrix}0&R\\R&0\end{pmatrix},
\qquad
\widehat J_s=\begin{pmatrix}0&R\\-R&0\end{pmatrix}.
\]
The unsigned \(\widehat J_0\) induces edge reversal but is anti-symplectic;
the signed \(\widehat J_s\) induces the same projective edge reversal and is
symplectic.

\begin{theorem}
\label{thm:20.1}
The finite pure \(\{C,J\}\) directed-edge geometry admits the signed
projective dihedral lift displayed above.  The unsigned reversal is
anti-symplectic, while the signed reversal is symplectic.  The finite tick word
\[
C,J,C,C,J
\]
visits all six directed edges and has a signed symplectic realization.
\end{theorem}

\begin{proof}
Direct multiplication gives
\[
\widehat C^T\Omega\widehat C=\Omega,
\qquad
\widehat J_0^T\Omega\widehat J_0=-\Omega,
\qquad
\widehat J_s^T\Omega\widehat J_s=\Omega,
\]
and
\[
\begin{aligned}
\widehat C^3&=I,
&
\widehat J_s^2&=-I,
&
\widehat J_s\widehat C\widehat J_s^{-1}&=\widehat C^{-1}.
\end{aligned}
\]
Hence the projective classes satisfy the dihedral relations in
\(\operatorname{PSp}(V,\omega)\).  On coordinate lines, \(\widehat C\) acts as
\((01\ 12\ 20)(10\ 02\ 21)\), and \(\widehat J_s\) acts as edge reversal.
Starting from \(01\), the tick word gives
\[
01\mapsto12\mapsto21\mapsto10\mapsto02\mapsto20.
\]
The signed-lift matrices and tick-word audit are recorded in the \texttt{L3/}
signed-symplectic bridge tables.
\end{proof}

The signed lift completes the post-selection passport of \(\RP\).


\section{Conclusion}\label{conclusion}

The normal-form landscape
\[
\Omega'\cong S^{21}
\]
is mapped by the associativity atlas and the \(\mathcal H\)-frontier.  The four
internal selector predicates
\[
\mathcal H_{acc}=0,
\qquad
H_{diag}=0,
\qquad
(r,r)^2=(r,r),
\qquad
\operatorname{AbsTrans}=\{C,J\}
\]
select the right-row projection magma \(\RP\).  After this selection, Part III
records the finite post-selection passport: operator envelope, term algebra,
endomorphism/intertwiner classification, row-contraction matrix bridge,
hidden-continuation operator identity, carrier-level compact companions, and
signed symplectic lift.

The manuscript therefore has the publication chain
\[
\begin{gathered}
\Omega'\cong S^{21}
\;\xrightarrow{\text{selector}}\;
\RP,
\\[1mm]
\RP
\;\xrightarrow{\text{linearization}}\;
\begin{matrix}
\text{finite operator, term,}\\
\text{and carrier passport}
\end{matrix}.
\end{gathered}
\]

\appendix
\clearpage

\phantomsection
\part*{Companion material}
\addcontentsline{toc}{part}{Companion material}


\section{Finite certificate protocol}\label{finite-certificate-protocol}

The finite support package is available in the public GitHub repository: 

\url{https://github.com/burundai-t/finite-row-graded-magmas-over-f3}.

The finite algebraic language follows the
standard binary-system, universal-algebra, design, quasigroup, and clone-theory
sources \cite{Birkhoff1935,Bruck1958,BurrisSankappanavar1981,Cohn1965,DenesKeedwell1974,ColbournDinitz2007,Pflugfelder1990,Post1941,Rosenberg1970,Szendrei1986,Lau2006}.
The certificate style is the standard finite-proof pattern: a human-readable
reduction is paired with reproducible finite artifacts
\cite{AppelHaken1977,AppelHakenKoch1977,Gonthier2008,Hales2005,HalesEtAl2017,McCune1997}.

\begin{longtable}[]{@{}p{0.24\linewidth}p{0.30\linewidth}p{0.38\linewidth}@{}}
\toprule\noalign{}
component & role & manuscript interface \\
\midrule\noalign{}
\endhead
\bottomrule\noalign{}
\endlastfoot
\texttt{L1/} & finite landscape and associativity reference checks & normal form, associativity atlas, endpoint loci, partition function \\
\texttt{L2/} & compact selection and pure \(C/J\) survivor checks & selector predicates, directed-edge absorption, uniqueness theorem \\
\texttt{L3/} & selected-\(\RP\) operator and bridge checks & operator envelope, \(T_2\), matrix bridge, compact companions, signed symplectic lift \\
\texttt{mathcal\_H/} & global \(\mathcal H\)-frontier bundle & local distance reduction, H1--H4 frontier obligations \\
repository interface & \path{README.md}; \path{SUPPORT_INDEX.txt} & reviewer commands and file-by-file index \\
\end{longtable}

The compact table families are:

\begin{longtable}[]{@{}p{0.25\linewidth}p{0.24\linewidth}p{0.43\linewidth}@{}}
\toprule\noalign{}
family & location & purpose \\
\midrule\noalign{}
\endhead
\bottomrule\noalign{}
\endlastfoot
L1 landscape atlas & \texttt{L1/tables/} & associativity distribution, endpoints, orbit representatives, controlled \(\mathcal H\) comparison \\
L2 selection & \texttt{L2/tables/} & access-zero, diagonal compression, pure \(C/J\) survivor separation \\
L3 linearization passport & \texttt{L3/tables/} & positive composition envelope, binary term algebra, matrix bridge, signed symplectic lift, fiber algebra \\
\(\mathcal H\) frontier & \texttt{mathcal\_H/} & H1 unrestricted range, H2 pure-frontier \(\mathrm{UNSAT}\), H3 witnesses, H4 spike classifier \\
\end{longtable}

\subsection[Global H-frontier certificate handles]{Global \(\mathcal H\)-frontier certificate handles}\label{global-h-frontier-certificate-handles}

The global \(\mathcal H\)-frontier theorem is checked by a separate
\texttt{mathcal\_H/} bundle.  At manuscript level the implication is
\[
\text{local distance reduction}
+
\text{H1--H4 certificate obligations}
\Longrightarrow
\text{global frontier theorem}.
\]

\begin{longtable}[]{@{}p{0.17\linewidth}p{0.30\linewidth}p{0.45\linewidth}@{}}
\toprule\noalign{}
handle & bundle location & theorem component \\
\midrule\noalign{}
\endhead
\bottomrule\noalign{}
\endlastfoot
HF-RED, HF-CERT-IF & bundle theorem interface & local distance reduction and local signature formulas \\
H1 & \texttt{h1/} & unrestricted full-landscape range \(rawH\in[-378,1217]\) \\
H2 & \texttt{h2/} & pure-frontier exclusion at \(rawH\ge1171\) \\
H3 & embedded witness verifier & \(\RP\) and \(\RC\) at \(rawH=1170\), \(N_-=0\) \\
H4 & \texttt{h4/} & twelve-point signed-cancellation spike at \(rawH=1217\) \\
\end{longtable}

The H2 obligation is the finite \(\mathrm{UNSAT}\) statement
\[
(h_q\ge0\ \forall q\in S^5)\wedge rawH\ge1171
\]
on the depth-9 live frontier:
\[
18540/18540\ \mathrm{UNSAT},
\qquad
\mathrm{SAT}=0,
\qquad
\mathrm{UNKNOWN}=0.
\]
The light static coverage verifier for the accepted H2 raw artifacts is:
\begin{verbatim}
cd mathcal_H/h2
PYTHONDONTWRITEBYTECODE=1 python3 verify_h2_coverage_light.py
\end{verbatim}

\subsection{Dependency summary}\label{dependency-summary}

\begin{longtable}[]{@{}p{0.36\linewidth}p{0.56\linewidth}@{}}
\toprule\noalign{}
result family & support role \\
\midrule\noalign{}
\endhead
\bottomrule\noalign{}
\endlastfoot
normal form and selector proof & hand finite arguments with L2 anchor checks \\
associativity atlas & block reduction plus L1 endpoint and partition-function tables \\
global \(\mathcal H\)-frontier & local distance reduction plus H1--H4 certificate bundle \\
operator envelope and row-fiber algebra & analytic row-owner reductions plus L3 multiplication tables \\
endomorphisms and intertwiners & finite row-fiber classification with linear intertwiner audit \\
binary term algebra & sector reduction plus realization and transition certificates \\
matrix, compact-companion, and symplectic bridges & carrier-level calculations supported by L3 bridge audits \\
\end{longtable}

\section[One-page passport of RP]{One-page passport of \(\RP\)}\label{one-page-passport-rp}

The selected point is
\[
\RP=(a,b,d)=(1,2,000).
\]
It has the finite economy and compression profile used by the selector:
\[
\mathcal H_{acc}=0,
\qquad
H_{diag}=0,
\qquad
(r,r)^2=(r,r),
\qquad
 a\ne b,
\qquad
 d=000.
\]
Thus its cross-row rule is column-blind, its same-point diagonal is compressed
and anchored, and its two cross-row anchors are nondegenerate.  The
Landauer/Shannon provenance of the first two predicates is zero hidden-column
access and zero diagonal entropy
\cite{Landauer1961,Bennett1973,Bennett1982,Shannon1948,CoverThomas2006}.

\begin{proposition}[Nondegenerate anchors]
\label{prop:10.1}
Among the nine column-blind rules \(h_{a,b}\), the condition \(a\ne b\) leaves
exactly six rules.
\end{proposition}

\begin{proof}
There are nine pairs \((a,b)\in S^2\), and the three pairs \((a,a)\) are
excluded.
\end{proof}

\begin{proposition}[Local associativity finalists]
\label{prop:10.3}
With \(d=000\), the six nondegenerate column-blind rules have
\[
\operatorname{Assoc}_{000}=219
\]
exactly at \((a,b)=(1,2)\) and \((2,1)\), namely \(\RP\) and \(\RC\).  The
other four nondegenerate rules have \(\operatorname{Assoc}_{000}=273\).
\end{proposition}

\begin{proof}
By Proposition~\ref{prop:4.3}, on the column-blind, \(d=000\) stratum,
\[
\operatorname{Assoc}_{000}(h_{a,b})
=
3\left(73+18[a=0]+18[b=0]+54[a=b]\right).
\]
On \(a\ne b\), this is minimized exactly when neither anchor is \(0\), giving
\((1,2)\) and \((2,1)\).
\end{proof}

The selected directed-edge feature is pure head-row absorption.

\begin{lemma}[Absorption passport]
\label{lem:12.1}
For \(\RP\),
\[
\operatorname{AbsTrans}_{\RP}=\{C,J\}.
\]
The row-complement comparison point \(\RC=(a,b,d)=(2,1,000)\) has
\[
\operatorname{AbsTrans}_{\RC}=\{C^{-1},C^{-1}J\}.
\]
\end{lemma}

\begin{proof}
Let \(x=(r,c)\in M^\times\).  For \(\RP\), the cross-row rule is
\(h(r_x,r_y)=r_y\), so absorption is \(r_y=c\).  The two absorbed neighbours
are \((c,r)\) and \((c,\overline{rc})\), namely \(J(x)\) and \(C(x)\).

For \(\RC\), the cross-row rule is \(h(r_x,r_y)=\overline{r_xr_y}\).  The
absorption equation \(\overline{rr_y}=c\) gives \(r_y=\overline{rc}\).  The two
absorbed neighbours are \((\overline{rc},r)\) and \((\overline{rc},c)\), namely
\(C^{-1}(x)\) and \(C^{-1}J(x)\).
\end{proof}

This one-page passport records \(\RP\) at the selection boundary: access
compression, diagonal compression, anchored self-action, nondegenerate anchors,
local associativity finalist value, and the head-row absorption pair
\(\{C,J\}\).  The adjacent row-complement finalist has the drifted absorption
pair \(\{C^{-1},C^{-1}J\}\).

\section{Theoretical-Informatic Reading of uniquely selected RP}
\label{structural-reading-selector}
\label{appendix-theoretical-informatic-reading}

The four selector predicates (together with the supporting structural properties of the selected magma RP) admit the following information-theoretic and structural reading:

\begin{center}
\begin{tabular}{p{0.48\textwidth} p{0.48\textwidth}}
\toprule
\textbf{Mathematical condition} & \textbf{Information-theoretic / structural reading} \\
\midrule
\(\mathcal{H}_{\mathrm{acc}} = 0\) 
& Do not read hidden columns; avoid Landauer-facing erasure of hidden information. \\[8pt]

\((a \ne b)\), for RP \(((a,b)=(1,2))\) 
& Non-triviality. The output depends on the right public direction; interaction does not collapse into an empty response. \\[8pt]

\(H_{\mathrm{diag}} = 0\) 
& Minimize Shannon entropy on the diagonal — do not distinguish diagonal types where distinction is unnecessary. \\[8pt]

\(\min\operatorname{Assoc}_{000}\) / low local associativity 
& The order of application has meaning (the sequence of operations carries information). \\[8pt]

pure \((C/J)\) and signed symplectic lift 
& Directed-edge structure is preserved without drift; at the linear level, phase-volume preservation appears. \\[8pt]

\(((r,c)^2 = (r,r))\) 
& Self-application returns to the element's own row-fixed point. \\
\bottomrule
\end{tabular}
\end{center}

These six readings can be distilled into two complementary axes of a single principle:

\begin{center}
\boxed{\text{Economize, but do not descend into triviality.}} \\[12pt]
\boxed{\text{Preserve structure and information — without unnecessary expenditure.}}
\end{center}

\section{Open bridges and extensions}\label{open-bridges-extensions}\label{open-problems-and-bridges}

The following fronts are subsequent work beyond the theorem package proved in
this manuscript.

\subsection[Functorial RP--gl(3) bridge]{Functorial \(\RP\)--\(\mathfrak{gl}(3)\) bridge}\label{functorial-rp-gl3-bridge}

The shared carrier
\[
k[M]\cong\operatorname{Mat}_3(k)
\]
supports the finite \(\RP\) multiplication and ordinary matrix multiplication.
The row-contraction bridge
\[
\mu_{cross}(A,B)=AJ_3B^\top
\]
aligns the cross-row branch with matrix rank-one contraction.  The next
problem is to identify a universal or functorial construction from the finite
magma to the matrix/Lie package.  The matrix and Lie background used in the
carrier-level bridge is standard \cite{HornJohnson2012,FultonHarris1991,Hall2015}.

\subsection[Symplectic development of pure C/J]{Symplectic development of pure \(C/J\)}\label{symplectic-development-of-pure-cj}

The pure directed-edge geometry has an explicit signed symplectic lift.  The
next layer is the associated continuous symplectic geometry: canonical
\((q,p)\)-charts, Hamiltonian paths, and coordinate-free separation between
head-row reversal and row-complement drift.  This is the natural continuation of
the finite lift in the language of symplectic geometry
\cite{GuilleminSternberg1990,McDuffSalamon2017}.

\subsection{Unified selection principle}\label{unified-selection-principle}

The uniqueness theorem uses four finite predicates: access-zero, diagonal
entropy zero, row-anchor idempotence, and pure \(C/J\).  A sharper formulation
would derive these predicates from one natural functional whose global minimizer
on \(\Omega'\) is the \(\RP\) representative.

\subsection{All-arity term algebra}\label{all-arity-term-algebra}

The binary term algebra satisfies
\[
T_2(\RP)\cong\{x,y\}\times\mathcal S_{21}\times\mathcal U_{15},
\qquad
|T_2(\RP)|=630.
\]
The all-arity clone remains to be classified.  The relevant background is the
classical clone-theoretic framework \cite{Post1941,Rosenberg1970,Szendrei1986,Lau2006}.

\subsection{Larger carriers}\label{larger-carriers}\label{larger-carriers-relaxed-ax2}

Natural extensions include analogues over larger finite fields or cyclic sets
and directed-edge \(C/J\) geometries beyond the triangle.  The standard
combinatorial-design and quasigroup references provide the ambient language for
these variants \cite{DenesKeedwell1974,ColbournDinitz2007,Pflugfelder1990}.

\subsection{Relaxing AX2}\label{relaxing-ax2}

The present landscape fixes the same-row Steiner rule.  The relaxed landscape
asks whether that same-row rule can itself be selected intrinsically.

\phantomsection
\section*{Note on AI assistance}
\addcontentsline{toc}{section}{Note on AI assistance}

This work was developed with assistance from large language models.  The author
specified the conceptual framework, constraints, and target properties of the
objects under study.  Within this framework, AI tools were used for
computational exploration, code drafting, and structuring intermediate results.
All mathematical definitions, statements, and theorems presented in the paper
were established and independently verified by the author.  The finite claims
are backed by exhaustive checks recorded in the accompanying support package.

\clearpage
\addcontentsline{toc}{section}{References}
\begin{thebibliography}{99}
\bibitem{AppelHaken1977} K. Appel and W. Haken, ``Every planar map is four colorable. Part I:
Discharging,'' \emph{Illinois Journal of Mathematics} \textbf{21}
(1977), no. 3, 429--490. doi:10.1215/ijm/1256049011.

\bibitem{AppelHakenKoch1977} K. Appel, W. Haken, and J. Koch, ``Every planar map is four colorable.
Part II: Reducibility,'' \emph{Illinois Journal of Mathematics}
\textbf{21} (1977), no. 3, 491--567. doi:10.1215/ijm/1256049012.

\bibitem{Bennett1973} C. H. Bennett, ``Logical reversibility of computation,'' \emph{IBM
Journal of Research and Development} \textbf{17} (1973), no. 6,
525--532. doi:10.1147/rd.176.0525.

\bibitem{Bennett1982} C. H. Bennett, ``The thermodynamics of computation---a review,''
\emph{International Journal of Theoretical Physics} \textbf{21} (1982),
no. 12, 905--940. doi:10.1007/BF02084158.

\bibitem{Birkhoff1935} G. Birkhoff, ``On the structure of abstract algebras,''
\emph{Proceedings of the Cambridge Philosophical Society} \textbf{31}
(1935), no. 4, 433--454. doi:10.1017/S0305004100013463.

\bibitem{Bruck1958} R. H. Bruck, \emph{A Survey of Binary Systems}, Ergebnisse der
Mathematik und ihrer Grenzgebiete, Neue Folge, Heft 20, Springer,
Berlin, 1958.

\bibitem{BurrisSankappanavar1981} S. Burris and H. P. Sankappanavar, \emph{A Course in Universal Algebra},
Graduate Texts in Mathematics 78, Springer-Verlag, New York, 1981.
doi:10.1007/978-1-4613-8130-3.

\bibitem{Cohn1965} P. M. Cohn, \emph{Universal Algebra}, Harper's Series in Modern
Mathematics, Harper \& Row, New York, 1965.

\bibitem{ColbournDinitz2007} C. J. Colbourn and J. H. Dinitz, eds., \emph{Handbook of Combinatorial
Designs}, 2nd ed., Chapman \& Hall/CRC, Boca Raton, 2007.

\bibitem{CoverThomas2006} T. M. Cover and J. A. Thomas, \emph{Elements of Information Theory}, 2nd
ed., Wiley-Interscience, Hoboken, 2006. doi:10.1002/047174882X.

\bibitem{DenesKeedwell1974} J. Dénes and A. D. Keedwell, \emph{Latin Squares and Their
Applications}, Academic Press, New York, 1974.

\bibitem{FultonHarris1991} W. Fulton and J. Harris, \emph{Representation Theory: A First Course},
Graduate Texts in Mathematics 129, Springer-Verlag, New York, 1991.
doi:10.1007/978-1-4612-0979-9.

\bibitem{Gonthier2008} G. Gonthier, ``Formal proof---the four-color theorem,'' \emph{Notices of
the American Mathematical Society} \textbf{55} (2008), no. 11,
1382--1393.

\bibitem{GuilleminSternberg1990} V. Guillemin and S. Sternberg, \emph{Symplectic Techniques in Physics},
Cambridge University Press, Cambridge, 1990.

\bibitem{Hales2005} T. C. Hales, ``A proof of the Kepler conjecture,'' \emph{Annals of
Mathematics} \textbf{162} (2005), no. 3, 1065--1185.
doi:10.4007/annals.2005.162.1065.

\bibitem{HalesEtAl2017} T. Hales, M. Adams, G. Bauer, T. D. Dang, J. Harrison, L. T. Hoang, C.
Kaliszyk, V. Magron, S. McLaughlin, T. T. Nguyen, Q. T. Nguyen, T.
Nipkow, S. Obua, J. Pleso, J. Rute, A. Solovyev, T. H. A. Ta, N. T.
Tran, T. D. Trieu, J. Urban, K. Vu, and R. Zumkeller, ``A formal proof
of the Kepler conjecture,'' \emph{Forum of Mathematics, Pi} \textbf{5}
(2017), e2. doi:10.1017/fmp.2017.1.

\bibitem{Hall2015} B. C. Hall, \emph{Lie Groups, Lie Algebras, and Representations: An
Elementary Introduction}, 2nd ed., Graduate Texts in Mathematics 222,
Springer, Cham, 2015. doi:10.1007/978-3-319-13467-3.

\bibitem{HornJohnson2012} R. A. Horn and C. R. Johnson, \emph{Matrix Analysis}, 2nd ed., Cambridge
University Press, Cambridge, 2012.

\bibitem{Landauer1961} R. Landauer, ``Irreversibility and heat generation in the computing
process,'' \emph{IBM Journal of Research and Development} \textbf{5}
(1961), no. 3, 183--191. doi:10.1147/rd.53.0183.

\bibitem{Lau2006} D. Lau, \emph{Function Algebras on Finite Sets: A Basic Course on
Many-Valued Logic and Clone Theory}, Springer Monographs in Mathematics,
Springer, Berlin, 2006. doi:10.1007/3-540-36023-9.

\bibitem{McCune1997} W. McCune, ``Solution of the Robbins problem,'' \emph{Journal of
Automated Reasoning} \textbf{19} (1997), no. 3, 263--276.
doi:10.1023/A:1005843212881.

\bibitem{McDuffSalamon2017} D. McDuff and D. Salamon, \emph{Introduction to Symplectic Topology},
3rd ed., Oxford Graduate Texts in Mathematics 27, Oxford University
Press, Oxford, 2017.

\bibitem{Pflugfelder1990} H. O. Pflugfelder, \emph{Quasigroups and Loops: Introduction}, Sigma
Series in Pure Mathematics 7, Heldermann Verlag, Berlin, 1990.

\bibitem{Post1941} E. L. Post, \emph{The Two-Valued Iterative Systems of Mathematical
Logic}, Annals of Mathematics Studies 5, Princeton University Press,
Princeton, 1941.

\bibitem{Rosenberg1970} I. G. Rosenberg, ``Über die funktionale Vollständigkeit in den
mehrwertigen Logiken,'' \emph{Rozpravy Československé Akademie Věd, Řada
matematických a přírodních věd} \textbf{80} (1970), no. 4, 3--93.

\bibitem{Shannon1948} C. E. Shannon, ``A mathematical theory of communication,'' \emph{Bell
System Technical Journal} \textbf{27} (1948), 379--423 and 623--656.
doi:10.1002/j.1538-7305.1948.tb01338.x;
doi:10.1002/j.1538-7305.1948.tb00917.x.

\bibitem{Szendrei1986} Á. Szendrei, \emph{Clones in Universal Algebra}, Séminaire de
mathématiques supérieures 99, Presses de l'Université de Montréal,
Montréal, 1986.
\end{thebibliography}

\end{document}
